\documentclass[11pt,a4paper]{article}
\usepackage[T1]{fontenc}
\usepackage[utf8]{inputenc}
\usepackage[margin=27mm,headheight=15pt]{geometry}
\usepackage{microtype}
\usepackage{amsmath,amsthm,mathtools}
\usepackage{newtxtext,newtxmath}
\usepackage{booktabs,array,longtable}
\usepackage[dvipsnames]{xcolor}
\usepackage{enumitem}
\usepackage{fancyhdr}
\usepackage[most]{tcolorbox}
\usepackage{hyperref}
\hypersetup{colorlinks=true,linkcolor=MidnightBlue,citecolor=MidnightBlue,
 urlcolor=MidnightBlue,pdftitle={Exact Precision and Algebraic Rigidity of Omnific Continued Fractions},
 pdfauthor={AI-assisted research draft},pdfsubject={Surreal numbers, Hahn fields, omnific continued fractions}}
\urlstyle{same}
\setlength{\emergencystretch}{2em}
\setlength{\parskip}{3pt}
\setlist{itemsep=3pt,topsep=5pt}
\pagestyle{fancy}
\fancyhf{}
\fancyhead[L]{\small Omnific continued fractions}
\fancyhead[R]{\small Precision and algebraic rigidity}
\fancyfoot[C]{\thepage}
\renewcommand{\headrulewidth}{0.3pt}
\numberwithin{equation}{section}
\newtheorem{theorem}{Theorem}[section]
\newtheorem{proposition}[theorem]{Proposition}
\newtheorem{lemma}[theorem]{Lemma}
\newtheorem{corollary}[theorem]{Corollary}
\theoremstyle{definition}
\newtheorem{definition}[theorem]{Definition}
\newtheorem{example}[theorem]{Example}
\newtheorem{problem}{Research question}
\theoremstyle{remark}
\newtheorem{remark}[theorem]{Remark}
\newcommand{\No}{\mathbf{No}}
\newcommand{\Oz}{\mathbf{Oz}}
\newcommand{\SC}{\mathbf{No}[i]}
\newcommand{\R}{\mathbb R}
\newcommand{\Q}{\mathbb Q}
\newcommand{\Z}{\mathbb Z}
\newcommand{\N}{\mathbb N}
\newcommand{\C}{\mathbb C}
\newcommand{\Ofin}{\mathcal O_{\mathrm{fin}}}% [write] the source wrote \mathcal O (macro \OO)
\newcommand{\mm}{\mathfrak m}
\newcommand{\F}{\mathcal F}
\newcommand{\alg}{\operatorname{Alg}}
\newcommand{\supp}{\operatorname{supp}}
\newcommand{\dgr}{\operatorname{deg}}
\newcommand{\fl}[1]{\left\lfloor#1\right\rfloor_{\Oz}}
\newcommand{\KG}{K_G}
\newcommand{\restr}{\mathbin{\upharpoonright}}
\newcommand{\codea}{\boldsymbol a}
\newcommand{\codeb}{\boldsymbol b}
\newcommand{\rpath}[1]{\nolinkurl{#1}}% a path in the collection, relative to docs/
\newcommand{\lbl}[1]{\nolinkurl{#1}}% a label of another report, cited by name
\newcommand{\wtag}{\textup{[write]}}% text added when the report joined the collection
\newtcolorbox{statusbox}{colback=black!3,colframe=black!35,boxrule=0.4pt,arc=1pt,
 left=9pt,right=9pt,top=7pt,bottom=7pt}
\newtcolorbox{resultbox}{colback=MidnightBlue!3,colframe=MidnightBlue!55,
 boxrule=0.5pt,arc=1pt,left=9pt,right=9pt,top=7pt,bottom=7pt,breakable}

\begin{document}
\hypersetup{pageanchor=false}
\begin{titlepage}
\centering
\vspace*{12mm}
{\small\scshape Surreal numbers \; / \; omnific arithmetic \; / \; Hahn fields\par}
\vspace{12mm}
{\LARGE\bfseries Exact Precision and Algebraic Rigidity\par}
\vspace{3mm}
{\LARGE\bfseries of Omnific Continued Fractions\par}
\vspace{7mm}
{\large Full Hahn-field realization, valuation ideals,\par
periodic codes, and surcomplex fibers\par}
\vspace{10mm}
{AI-assisted research draft\par}
{23 September 2026\par}
\vspace{10mm}
\begin{minipage}{0.93\textwidth}
\begin{abstract}
An ordinary infinite continued fraction with omnific integer digits does not
usually determine one surreal number. Its indeterminacy was already studied
by Roggeman. We give a valuation-theoretic refinement and use it to derive
precise realization and algebraicity criteria. If $q_n$ are the continuant
denominators and $d_n=\dgr q_n$, the complete fiber of a digit sequence is
exactly a translate of the ideal of errors $\epsilon$ satisfying
$v(\epsilon)>2d_n$ for every ordinary $n$. A uniform two-step buffer estimate
makes this description independent of the chosen realization.

Every admissible sequence in a full real-coefficient Hahn field has a
realization in that same field; it has exactly one precisely when the
cumulative denominator degrees are cofinal in the value group. We classify
when the precision ideal is prime and give an exact two-condition criterion
for the fiber to contain at most one element algebraic over its digit field.
For eventually periodic codes an algebraic realization always exists, so
this becomes a sharp one-versus-countably-many dichotomy. Each full surreal
fiber nevertheless contains algebraically independent families of every
set cardinality over every prescribed set-sized coefficient field.
Coordinatewise surcomplex fibers and their changes under fractional linear
maps complete the picture. The proofs are written in full, with finite
symbolic checks and an explicit research agenda.
\end{abstract}
\vspace{3mm}
\begin{statusbox}
\small\textbf{Status and priority.} This is an unrefereed, AI-assisted research
draft, not a certified resolution of a named longstanding open problem.
Basic surreal continued-fraction nonuniqueness is \emph{not} claimed as new.
The proposed contributions are the sharpened precision, workspace, and
algebraic-fiber classifications. Their novelty has not been independently
certified. The accompanying finite tests are not a Lean verification or a
proof of the infinite-support assertions.

\smallskip
\wtag{} Single-source research report of the collection, built from
manuscript~09 of batch~33 and placed in commit \texttt{aa9c891}.
Independent proof review and formalization are pending. Labels carry the
prefix \texttt{ocf:}. Conventions, the place in the collection, the
non-claims and the provenance are in
Sections~\ref{ocf:sec:conventions}, \ref{ocf:sec:collection},
\ref{ocf:sec:nonclaims} and Appendix~\ref{ocf:app:report}; added text is
marked \wtag.
\end{statusbox}
\end{minipage}
\vfill
{\small Repository comparison pinned to \texttt{958b5c486581}\par}
\end{titlepage}
\pagenumbering{roman}
\hypersetup{pageanchor=true}

\begingroup
\small
\setlength{\parskip}{0pt}
\tableofcontents
\endgroup
\newpage
\pagenumbering{arabic}

\section{The question and the contribution}
\label{ocf:sec:intro}

Over the real numbers, an infinite regular continued-fraction expansion is a
representation of a single irrational number. Over the surreal numbers,
there are two different questions: which points produce a specified digit
sequence, and which point should be selected as its distinguished value?
The first is an arithmetic question about the floor algorithm. The second
requires an additional selection principle. Confusing them makes a periodic
code look more rigid than it is.

This article studies the first question, with digits in Conway's omnific
integer ring $\Oz$. The sequence length is exactly $\omega$: all indices
$n,k,m$ used in the continued-fraction algorithm are \emph{ordinary finite}
natural numbers. Neither a limit-stage continuation of the algorithm nor a
Gaussian nearest-integer algorithm is silently included.

\subsection{What was already known}
\label{ocf:sub:known}

Roggeman's \emph{Continued fractions and surreal numbers}
\cite{Roggeman} already develops the omnific floor algorithm, associated
classes of points with the same infinite code, and periodic fractions. In
particular, his Section~3A characterizes associated classes by convergent
inequalities; Section~3H studies their size and structure; Section~3I
introduces a generalized continuation. Thus nonuniqueness and a
proper-class fiber are antecedents, not discoveries of the present article.
The relevant scanned pages were inspected directly.

The continuant identities and the real periodicity theorem are classical;
\cite{Krishnan,Eberl} give accessible treatments, with \cite{Eberl} providing
a formalized real theory. Our finite identities are rederived because the
infinite digits and non-Archimedean inequalities require careful hypotheses.
Conway's and Gonshor's normal-form and real-closedness theorems are the
surreal foundations used here \cite{Conway,Gonshor}.

The repository \cite{SurrealRepo} supplies the setting and an implementation
roadmap. The inspected root README and report catalogue include the actual
omnific floor, degree, and normal-form infrastructure. No continued-fraction
report was identified in that catalogue. This was a focused comparison, not
a declaration-by-declaration or whole-history audit. No unreviewed advanced
repository theorem is required below.

\wtag{} This subsection is the source's, written at its pin
\texttt{958b5c4}. Its statement about the collection is still true at the
placement commit \texttt{aa9c891}: no report treats surreal or omnific
continued fractions, and Roggeman is cited nowhere else. The only continued
fractions elsewhere are ordinary real ones (a rotation number in
\rpath{surcomplex/dynamics-and-normal-forms}, \lbl{dyn:prop:cf}) and the
tridiagonal determinant continuants of
\rpath{surreal/matrix-scaling-at-surreal-scales}
(\lbl{scale:eq:continuantval}). The omnific floor of
Lemma~\ref{ocf:lem:floor} is the collection's integer-part theorem, proved
there and formalized in Lean. The relation to the collection, with labels,
is in Section~\ref{ocf:sec:collection}.

\subsection{The central classification}

Let $\codea=(a_0,a_1,\ldots)$, where $a_0\in\Oz$ and $a_n\in\Oz$,
$a_n\ge1$ for $n\ge1$. Write $\F(\codea)$ for all surreals whose
nonterminating floor algorithm produces this code. Define the usual
continuant denominators by
\[
 q_{-1}=0,\qquad q_0=1,\qquad
 q_n=a_nq_{n-1}+q_{n-2}\quad(n\ge1),
\]
and put
\begin{equation}\label{ocf:eq:summarydegrees}
 d_0=0,\qquad d_n=\sum_{j=1}^{n}\dgr a_j=\dgr q_n.
\end{equation}
Here $\dgr$ denotes the greatest \emph{growth exponent} in Conway normal
form; the natural valuation has the opposite sign.

\begin{resultbox}
For every $x\in\F(\codea)$,
\[
 \F(\codea)=x+I_{\codea},\qquad
 I_{\codea}=\{0\}\cup
 \{\epsilon\ne0:v(\epsilon)>2d_n\text{ for every }n\}.
\]
The sequence records a coset of an explicit valuation ideal, rather than a
single unrestricted surreal number. The ideal depends only on the degrees
of the digits after $a_0$, not on their other coefficients.
\end{resultbox}

The exact formula is proved in Theorem~\ref{ocf:thm:fiber}, using a buffer from
both cylinder endpoints. It refines the earlier associated-class viewpoint
into a form usable for Hahn-field descent and algebraic rigidity.

\subsection{Consequences that go beyond basic nonuniqueness}

Theorem~\ref{ocf:thm:hahn} proves existence in every \emph{full}
$\R$-coefficient Hahn field containing the digits, not merely existence
somewhere in $\No$. Its uniqueness criterion is exactly cofinality of
$(d_n)$ in the value group. The existence proof is not a completeness
assertion about arbitrary non-Archimedean ordered fields.

For the countable digit field $K_{\codea}=\Q(a_0,a_1,\ldots)$, set
$D_0=\dgr a_0$ when $a_0\ne0$ and $D_0=0$ otherwise. Two simple conditions
control algebraic rigidity:
\begin{align}
 &\text{for every }n\text{ there is }m\text{ with }d_m\ge2d_n,
    \tag{P}\label{ocf:cond:P}\\
 &\text{there is }n\text{ with }D_0\le2d_n.
    \tag{V}\label{ocf:cond:V}
\end{align}
When both hold, the fiber has at most one element algebraic over
$K_{\codea}$. If either fails and one algebraic realization exists, there are
countably infinitely many, all obtainable in the same finite field
extension by rational translations. This is
Theorem~\ref{ocf:thm:algebraic}. For eventually periodic codes, existence of an
algebraic realization turns this conditional statement into an exact
dichotomy (Theorem~\ref{ocf:thm:periodiccriterion}).

The prime-ideal criterion, the algebraic dichotomy, and the universal
Hahn-workspace realization theorem are the main proposed research
contributions. We do not assert that a literature search establishes their
priority. Their mathematical content is specified by the statements and
proofs, not by a label such as ``breakthrough.''

\subsection{Notation and conventions in this report}
\label{ocf:sec:conventions}

\wtag{} This subsection and the next were written when the manuscript
joined the collection. One symbol was renamed: the ring of finite surreals,
written $\mathcal O$ in the source, is written $\Ofin$ here, as in the
omnific Diophantine report (\rpath{surreal/omnific-diophantine-geometry}),
where $\operatorname{O}$ is an orthogonal group and $\mathcal O_\Bbbk(\Gamma)$ a
Hahn valuation ring. Its complex counterpart, written $\mathcal
O_\C=\mathcal O+i\mathcal O$ in the source, is written $\Ofin[i]$ in
Theorem~\ref{ocf:thm:complex}. No normalization changed. All other symbols
are the source's; their meanings are fixed below for the whole report, and
collisions are disclosed rather than removed. Paths such as
\rpath{surreal/omnific-diophantine-geometry} are relative to
\rpath{docs/}, and labels of other reports are cited by name.
\begin{itemize}[leftmargin=*,itemsep=3pt]
\item \emph{Exponents and valuation.} The report uses Conway's decreasing
large-$\omega$ convention $x=\sum r_\gamma\omega^\gamma$ throughout; no
$t^\gamma=\omega^{-\gamma}$ convention occurs. $\dgr x$ is the greatest
growth exponent and $v=-\dgr$, as in the collection's notation guide
(\rpath{NOTATION.md}) and as the Diophantine report's
$\operatorname{deg}_\omega$, which sets $\operatorname{deg}_\omega0=-\infty$;
here $\dgr0$ is never used. A \emph{large} valuation means a \emph{small}
element: the tempting reading of $v(e)>2d_n$ as ``$e$ is large'' is false;
it says $\lvert e\rvert\prec q_n^{-2}$.
\item \emph{Degrees.} $d_n=\dgr q_n=\sum_{j=1}^n\dgr a_j$; the first digit
never contributes. $D_0$ is $\dgr a_0$, or $0$ when $a_0=0$: a bookkeeping
value, not a degree of zero. $D_*$ and $E$ in \eqref{ocf:eq:periodscales}
are the largest displayed degree and the degree gained per period.
\item \emph{Rings and ideals.} $\Oz$ is the omnific ring, $\Ofin$ the finite
surreals and $\mm$ the infinitesimals. The precision ideal $I_{\codea}$ is a
convex ideal of the valuation ring $\Ofin$ contained in $\mm$. It is
\emph{not} an ideal of $\Oz$ (it contains no nonzero omnific integer) or of
$\No$, not the purely infinite ideal $\Pi$ of other reports, and not one of
the ideals $I_\kappa$ of $\Oz$ in \rpath{surreal/set-sized-quotients-of-omnific-integers}.
``Prime'' always refers to $\Ofin$ or $\Ofin\cap K_G$.
\item \emph{Fields.} $K_G=\R((\omega^G))$ is the \emph{full} Hahn field on a
set-sized additive subgroup $G\subseteq\No$, which need be neither
divisible nor convex; it is the field $\mathcal K_\R(G)$ of
\rpath{surreal/independent-surreal-copies}. A ``set-sized Hahn workspace'' is
such a $K_G$ containing the data; the workspaces of the foundations report
(\lbl{found:thm:workspace}) are of this kind with $G$ divisible, and no
divisibility is needed here. $K_{\codea}=\Q(a_0,a_1,\ldots)$ is the
countable field generated by the digits, and $K=K_{\codea}$ in
Section~\ref{ocf:sec:periodic}. The tempting reading of $K_{\codea}$ as a
Hahn field $K_G$ is false; $K_\R$ in Section~\ref{ocf:sec:examples} is $K_G$
at $G=\R$.
\item \emph{Fibre and Fibonacci.} $\F(\codea)$ (calligraphic) is the fibre
of a code; $F_k$ (italic) are the Fibonacci numbers of \eqref{ocf:eq:fib}
and of the proof of Theorem~\ref{ocf:thm:fiber}. $A_{\codea}$ is the
algebraic part of a fibre (Section~\ref{ocf:sec:algebraic}), while $A,B,C,D$
are the entries of the period matrix $M$ in \eqref{ocf:eq:periodmatrix}.
\item \emph{Letters with several local meanings.}
$D$ is an entry of $M$ in Section~\ref{ocf:sec:periodic} and the lower set
$\bigcup_n[0,2d_n]$ in the proof of Theorem~\ref{ocf:thm:prime}; neither is
$D_0$ or $D_*$.
$E$ is the period degree in Section~\ref{ocf:sec:periodic} (and in
Section~\ref{ocf:sec:collection}), but a set-sized subfield in
Lemma~\ref{ocf:lem:algerror}, Section~\ref{ocf:sec:independence} and
Section~\ref{ocf:sec:complex}; $E_t$ is a partial period sum.
$t$ is the variable of $\Phi_n$ and the complete quotient $x_{n+1}$ in
Sections~\ref{ocf:sec:cylinders}--\ref{ocf:sec:fibers}, the period-fixed
root in Theorem~\ref{ocf:thm:periodroot}, and an ordinary index
$0\le t<\ell$ in \eqref{ocf:eq:perioddegree} and the proofs of
Theorems~\ref{ocf:thm:periodiccriterion} and~\ref{ocf:thm:orderunit}:
$E_t$ is not $E$ evaluated at the root.
$q_n$ are continuant denominators, while a bare $q\in\Q$ is the rational
multiplier in \eqref{ocf:eq:algtranslations}, in the proof of
Theorem~\ref{ocf:thm:periodiccriterion} and in Section~\ref{ocf:sec:examples}:
$\xi+q\delta$ is not $\xi+q_n\delta$.
$r_n$ are convergents, $r$ is a real constant coefficient in
Lemma~\ref{ocf:lem:floor} and the prefix length in
\eqref{ocf:eq:periodiccode}; $s_n$ is the other cylinder endpoint.
$S$ is a support in \eqref{ocf:eq:normal}, a set of distances in the proof
of Proposition~\ref{ocf:prop:topology}, and $\sum_{j<r}\dgr a_j$ in
Section~\ref{ocf:sec:periodic}.
$C_n$ are cylinders, $C$ an entry of $M$, $C_k$ Catalan numbers in
\eqref{ocf:eq:rho}, and $\C$ the complex numbers.
$M$ is the period matrix and a bound on total degrees in the proof of
Lemma~\ref{ocf:lem:digitfield}; $N$ is an index in
Theorem~\ref{ocf:thm:hahn} and Lemma~\ref{ocf:lem:digitfield} and an
ordinary multiplier in \eqref{ocf:eq:periodcriterion}; $m$ is an ordinary
integer and $\mm$ the ideal.
$a,b$ are the digits $a_{n+1},a_{n+2}$ in the proof of
Lemma~\ref{ocf:lem:buffer} and elements of $\Ofin$ in the proof of
Theorem~\ref{ocf:thm:complex}; $a$ is the floor value in
Lemma~\ref{ocf:lem:floor}; $b_j$ are period digits and $b$ a digit of
largest degree in Section~\ref{ocf:sec:periodic}; $\codeb$ is a second code
in Section~\ref{ocf:sec:complex}; and $a,b$,
with $c,d$, the complex coefficients of the map in
Proposition~\ref{ocf:prop:mobius}, where $d$ is not a degree $d_n$.
$e$ is an error, except that it is an order unit in
Theorem~\ref{ocf:thm:orderunit}. $\delta$ is a nonzero element of
$I_{\codea}\cap K_{\codea}$ in Sections~\ref{ocf:sec:algebraic}
and~\ref{ocf:sec:periodic}, a positive bound in the proof of
Proposition~\ref{ocf:prop:topology}, and $\delta_\alpha$ are the
perturbations of Section~\ref{ocf:sec:independence}. $\gamma$ is an
exponent in \eqref{ocf:eq:normal}, a bound in the proof of
Corollary~\ref{ocf:cor:neverunique}, and $\gamma_n$ a prescribed cut in
Proposition~\ref{ocf:prop:prescribed}. $u,w$ are complete quotients in the
proof of Lemma~\ref{ocf:lem:buffer}, factors in the proof of
Corollary~\ref{ocf:cor:radical}, and coordinates in $v_\C(u+iw)$. $\rho$ is a real value in the proof of
Theorem~\ref{ocf:thm:hahn} and the fixed point \eqref{ocf:eq:rho}.
$\kappa$ is a set cardinal; it is neither the scale $\kappa=v(q)$ of
\rpath{surcomplex/expanding-polynomial-dynamics} nor a critical point.
\item \emph{Maps.} $\Phi_n$ of \eqref{ocf:eq:Phi} is a finite prefix map. It
is not the monomial lift $\Phi_h$ of the independent-copies report and not a
Taylor automorphism $\Phi_d$. $T$ is the period map. The support
projection $\pi_G$ in the proof of Theorem~\ref{ocf:thm:hahn} keeps the
terms with exponents in $G$: it is the bounded-support report's coset
projection at the zero coset (\lbl{bst:eq:projection}), additive and not
multiplicative.
\item \emph{Topologies.} Proposition~\ref{ocf:prop:topology} concerns the
order topology of the class $\No$, Corollary~\ref{ocf:cor:convergence} that of the set $K_G$. Continuants never
converge in the first and do converge in the second in the unique case;
the two statements do not conflict.
\item All indices $n,k,m,j$ of the algorithm are ordinary natural numbers.
The spelling ``fiber'' is the source's.
\end{itemize}

\subsection{Place in the collection}\label{ocf:sec:collection}

\wtag{} The source compared itself with the catalogue at its pin
(Section~\ref{ocf:sub:known}; Appendix~\ref{ocf:app:pinned}). Checked
against the collection at the placement commit, the relations are these.
No named question of any report is answered, and no report is contradicted.

\paragraph{The integer part.} Lemma~\ref{ocf:lem:floor} is the
integer-part theorem \lbl{odg:thm:floor} of the omnific Diophantine report,
with the same formula \lbl{odg:eq:floor}; its example
$\lfloor-\omega^{-1}\rfloor_{\Oz}=-1$ is the source's boundary case
$\fl{-e}=-1$. The collection proves it also in
the definable-surreals report (\lbl{dsn:prop:floor}) and the
computer-algebra report (\lbl{cas:eq-floor}), and prints it once in the
omnific-notations report (\lbl{onot:eq:floor}), which lists these
re-proofs; all three are in \rpath{foundations-and-computation/}. Here it is printed
once, with the source's proof, as one more independent re-derivation; no
second route is kept. It is the only ingredient of this report that is
formalized: the ledger \rpath{FORMALIZATION.md} marks \lbl{odg:thm:floor}
and \lbl{odg:eq:floor} \textbf{Proved} in
\rpath{../Surreal/Foundations/OmnificFloor.lean} (among others \texttt{omnificFloor\_spec} and
\texttt{omnificFloor\_of\_negative\_infinitesimal}). The label
\texttt{ocf:lem:floor} itself has no Lean mapping.

\paragraph{Itinerary fibres at an order unit.}
\rpath{surcomplex/expanding-polynomial-dynamics} (prefix \texttt{epd:})
codes the integral points of an expanding polynomial over a Hahn field by
words, and proves that each word's fibre is a centre plus the ideal
$I_\kappa=\{h:v(h)>n\kappa\text{ for all }n\}$ (\lbl{epd:lem:ideal},
\lbl{epd:thm:fibers}), so that a word determines one point exactly when the
scale $\kappa$ is an order unit (\lbl{epd:def:orderunit}). The exact fibre
theorem~\ref{ocf:thm:fiber} has the same shape for a different coding, the
omnific floor algorithm, with an arbitrary countable degree cut in place of
the multiples of one scale. For an eventually periodic code with $E>0$ and
$\sum_{j<r}\dgr a_j\le NE$ for some ordinary $N$, one has
$I_{\codea}=\{0\}\cup\{e:v(e)>kE\text{ for every }k\}$, which is that ideal
at $\kappa=E$, and Theorem~\ref{ocf:thm:orderunit} is the same order-unit
threshold. The two proofs are independent; neither report uses the other.
The separation of countable cofinality from an order unit
(Corollary~\ref{ocf:cor:cofinality}, Theorem~\ref{ocf:thm:orderunit},
Section~\ref{ocf:sec:examples}) is also the dividing line of the
ideal-theoretic trichotomy of
\rpath{surcomplex/entire-functions-at-arbitrary-rank}
(\lbl{ent:def:order-unit}); its group without an order unit is of the kind
in \lbl{epd:ex:no-order-unit}.

\paragraph{Support cosets.} The independence argument of
Theorem~\ref{ocf:thm:independence}, in which new exponents are rationally
independent modulo the span $H$ of all given supports so that distinct
monomials occupy disjoint support cosets, is the standard mechanism of
\rpath{surreal/independent-surreal-copies} (coset slices,
\lbl{isc:lem:slice}; linear disjointness of full Hahn fields over the field
of the intersected group, \lbl{isc:thm:sliceddisjoint}) and of
\rpath{surreal/transcendence-over-bounded-support} (coset projections,
\lbl{bst:eq:projection}, \lbl{bst:lem:projection}). It is credited to that
mechanism; the source's direct proof is printed. The algebraic independence
of the $\delta_\alpha$ over $\R((\omega^H))$ can also be read off from
\lbl{isc:thm:sliceddisjoint}, since $H$ meets the rational span of the new
exponents only in $0$.

\paragraph{Unrelated occurrences.} The continued fractions of
\rpath{surcomplex/dynamics-and-normal-forms} (\lbl{dyn:prop:cf},
\lbl{dyn:thm:nonbrjuno}) are ordinary real expansions of a rotation
number; the continuants of \rpath{surreal/matrix-scaling-at-surreal-scales}
(\lbl{scale:eq:continuantval}) are tridiagonal determinants. Their
valuation is computed by a leading-term induction like that of
Lemma~\ref{ocf:lem:degrees}, but their recurrence has a minus sign, so
strict dominance of one term replaces the positivity used here.

\paragraph{Formalization.} No \texttt{ocf:} label has a Lean mapping. The
collection's Lean library has no continuant or continued-fraction
declaration. Mathlib's continued fractions (\texttt{GenContFract.of})
require a floor to $\Z$ (\texttt{FloorRing}), which the omnific floor is
not, so they do not apply to $\No$ as they stand. The formalization route of
Section~\ref{ocf:sub:formal} is a proposal.

\paragraph{Status of the research questions.} All twelve questions of
Section~\ref{ocf:sec:questions} are new to the collection and open. None
is answered by another report. Question~\ref{ocf:q:formal} can start from
the proved floor. Question~\ref{ocf:q:surcomplex} is not addressed by
any surcomplex report, and Question~\ref{ocf:q:birthday} is not decided by
the birthday reports (\rpath{surreal/gonshor-product-birthdays},
\rpath{surreal/gonshor-laurent-birthdays}).

\section{Surreal and Hahn conventions}
\label{ocf:sec:prelim}

\subsection{Normal forms, degree, and valuation}

We use the following standard normal-form facts as inputs
\cite{Gonshor}. Every nonzero surreal has a unique expansion
\begin{equation}\label{ocf:eq:normal}
 x=\sum_{\gamma\in S}r_\gamma\omega^\gamma,
 \qquad r_\gamma\in\R\setminus\{0\},
\end{equation}
where $S$ is a reverse well-ordered \emph{set} of surreals. Every such
normal-form series exists. Its greatest exponent and leading coefficient
determine its growth and sign. Addition and multiplication are the Hahn
operations, and $\omega^{\gamma+\delta}=\omega^\gamma\omega^\delta$.
We write
\[
 \dgr x=\max S,\qquad v(x)=-\dgr x,\qquad v(0)=+\infty.
\]
In particular $v(xy)=v(x)+v(y)$ for nonzero $x,y$, and
$v(x+y)\ge\min\{v(x),v(y)\}$, with equality when the two valuations differ.
For positive $x,y$, leading coefficients do not cancel, so
\begin{equation}\label{ocf:eq:positivedegree}
 \dgr(x+y)=\max\{\dgr x,\dgr y\}.
\end{equation}
The symbol $\dgr0$ is not used in these identities. The convention
$D_0=0$ for a zero first digit is a separate bookkeeping definition.

Let $\Ofin$ be the ring of finite surreals and $\mm$ its infinitesimal ideal:
\[
 \Ofin=\{0\}\cup\{x:v(x)\ge0\},\qquad
 \mm=\{0\}\cup\{x:v(x)>0\}.
\]
The residue field is $\R$. For $b>0$ we use
\begin{equation}\label{ocf:eq:little}
 |e|\prec b\quad\Longleftrightarrow\quad
 |e|<b/m\ \text{for every ordinary integer }m\ge1.
\end{equation}
For nonzero $e$, this is equivalent to $v(e)>v(b)$. Indeed, a leading
valuation equality makes $|e|/b$ have a positive real leading coefficient,
which cannot be smaller than every positive real reciprocal.

\subsection{The omnific floor}

\begin{definition}\label{ocf:def:Oz}
The omnific ring $\Oz$ consists of the normal forms with growth exponents
in $\No_{\ge0}$ and with ordinary integer coefficient at exponent zero.
The purely infinite coefficients may be arbitrary real numbers.
\end{definition}

A nonzero omnific integer has nonnegative degree. An omnific integer of
degree zero is an ordinary nonzero integer. The ring is discretely ordered:
its least positive element is $1$. These observations follow immediately
from the leading term and the integer constant coefficient.

\begin{lemma}[Exact floor, including the boundary correction]
\label{ocf:lem:floor}
Each $x\in\No$ has a unique $a\in\Oz$ with $a\le x<a+1$. Write
$x=x_{>0}+r+e$, where $x_{>0}$ is its positive-growth part, $r\in\R$, and
$e\in\mm$. If $r\notin\Z$, then
$a=x_{>0}+\lfloor r\rfloor$. If $r\in\Z$, then
\[
 a=\begin{cases}x_{>0}+r,&e\ge0,\\
                 x_{>0}+r-1,&e<0.
   \end{cases}
\]
\end{lemma}
\begin{proof}
In the nonintegral case the real fractional part of $r$ is strictly between
zero and one, and adding an infinitesimal preserves both inequalities. In
the integral case the displayed correction gives the same bounds. For
uniqueness, two distinct omnific integers differ in absolute value by at
least one and therefore cannot both lie in $(x-1,x]$.
\end{proof}

We denote this floor by $\fl{x}$. In particular, if $0<e\in\mm$, then
$\fl{-e}=-1$, not zero. The algorithm below needs the actual floor, not
constant-term rounding without this correction.

\wtag{} Lemma~\ref{ocf:lem:floor} is the collection's integer-part
theorem \lbl{odg:thm:floor} (formula \lbl{odg:eq:floor}), also
\lbl{dsn:prop:floor}, \lbl{cas:eq-floor} and \lbl{onot:eq:floor}. It is
printed once here, with the source's proof, as an independent
re-derivation. The ledger marks \lbl{odg:thm:floor} \textbf{Proved} in Lean
(\rpath{../Surreal/Foundations/OmnificFloor.lean}), including the
negative-infinitesimal correction $\fl{-e}=-1$
(Section~\ref{ocf:sec:collection}).

\subsection{Full Hahn fields and class size}

For a set-sized additive subgroup $G\subseteq\No$, let
\begin{equation}\label{ocf:eq:KG}
 K_G=\R((\omega^G))
   =\{x\in\No:\supp(x)\subseteq G\},\qquad
 \Oz_G=K_G\cap\Oz.
\end{equation}
The support in this formula is the set of growth exponents. By the normal
form theorem $K_G$ is a set-sized ordered field, with value group $G$.
It is the \emph{full} Hahn field: arbitrary allowed well-ordered valuation
supports are admitted, not only finite or grid-based supports. Formula
\eqref{ocf:eq:KG} includes $G=\{0\}$, when $K_G=\R$.

The floor in Lemma~\ref{ocf:lem:floor} preserves $K_G$, since it deletes terms
and changes only the constant coefficient. Consequently the continued
fraction of any element of $K_G$, when defined, has digits in $\Oz_G$.

Class statements can be read in a standard conservative class framework,
or as assertions about the definable class of surreal sign sequences over
ZFC. Every support, sequence, field generated by specified data, and cut
used in a construction below is a set. No large-cardinal assumption,
class-length recursion for this algorithm, or proper-class choice of a
transcendence basis is used.

\section{Finite continuants and open cylinders}
\label{ocf:sec:cylinders}

\begin{definition}[Digits and complete quotients]\label{ocf:def:digits}
For $x_0=x$, define recursively
\[
 a_n=\fl{x_n},\qquad x_{n+1}=\frac1{x_n-a_n}
\]
provided $x_n\ne a_n$. A sequence $\codea=(a_n)_{n\ge0}$ is
\emph{admissible} if $a_0\in\Oz$ and $a_n\in\Oz$, $a_n\ge1$ for every
$n\ge1$. Its fiber $\F(\codea)$ consists of inputs for which all complete
quotients exist and the digit sequence is exactly $\codea$.
\end{definition}

The remainder belongs to $(0,1)$ whenever the algorithm continues. Hence
$x_{n+1}>1$ and its floor is at least one. Finite terminating expansions
are not included in an infinite fiber.

For an admissible code define, including the initial indices,
\begin{align}\label{ocf:eq:continuants}
 p_{-2}=0,\quad p_{-1}=1,&\qquad
 p_n=a_np_{n-1}+p_{n-2},\\
 q_{-2}=1,\quad q_{-1}=0,&\qquad
 q_n=a_nq_{n-1}+q_{n-2}.\nonumber
\end{align}
Then $p_0=a_0$ and $q_0=1$. All entries belong to $\Oz$.

\begin{lemma}[Matrix identities]\label{ocf:lem:matrix}
For every $n\ge0$,
\begin{align}
 \prod_{j=0}^{n}\begin{pmatrix}a_j&1\\1&0\end{pmatrix}
  &=\begin{pmatrix}p_n&p_{n-1}\\q_n&q_{n-1}\end{pmatrix},\label{ocf:eq:matrix}\\
 p_nq_{n-1}-p_{n-1}q_n&=(-1)^{n+1}.\label{ocf:eq:det}
\end{align}
Moreover $q_n>0$ and $q_n\ge q_{n-1}$, and, for the Fibonacci numbers
$F_1=F_2=1$,
\begin{equation}\label{ocf:eq:fib}
 q_{n+k}\ge F_{k+1}q_n\qquad(n,k\ge0).
\end{equation}
\end{lemma}
\begin{proof}
Multiplication of the last matrix gives the recursions in
\eqref{ocf:eq:continuants}; each factor has determinant $-1$. For the positive
and monotonic assertions start with $q_{-1}=0,q_0=1$ and use $a_j\ge1$.
The recurrence gives $q_{j+2}\ge q_{j+1}+q_j$. Induction on $k$, starting
with $k=0,1$, proves \eqref{ocf:eq:fib}. No Archimedean hypothesis enters.
\end{proof}

Define
\begin{equation}\label{ocf:eq:Phi}
 \Phi_n(t)=\frac{p_nt+p_{n-1}}{q_nt+q_{n-1}},\qquad
 r_n=\frac{p_n}{q_n},\qquad
 s_n=\frac{p_n+p_{n-1}}{q_n+q_{n-1}}.
\end{equation}
The map $\Phi_n$ is the finite composition
$a_0+1/(a_1+1/(\cdots+1/(a_n+1/t)))$.

\begin{proposition}[Cylinder description]\label{ocf:prop:cylinder}
The points having the first $n+1$ digits $a_0,\ldots,a_n$ and a further
complete quotient form the open interval
\[
 C_n=\Phi_n((1,\infty))
    =(\min(r_n,s_n),\max(r_n,s_n)).
\]
These intervals are nested, and
\begin{equation}\label{ocf:eq:width}
 |C_n|=\frac1{q_n(q_n+q_{n-1})}\le\frac1{q_n^2}.
\end{equation}
If $x=\Phi_n(t)$ with $t>1$, its distances to the two endpoints are
\begin{align}
 |x-r_n|&=\frac1{q_n(q_nt+q_{n-1})},\label{ocf:eq:distR}\\
 |x-s_n|&=\frac{t-1}{(q_n+q_{n-1})(q_nt+q_{n-1})}.\label{ocf:eq:distS}
\end{align}
Finally,
\begin{equation}\label{ocf:eq:intersection}
 \F(\codea)=\bigcap_{n\ge0}C_n.
\end{equation}
\end{proposition}
\begin{proof}
Each map $u\mapsto a+1/u$ takes $u>1$ into $(a,a+1)$, so its floor is
$a$. Iterating proves the digit interpretation and nesting. Subtracting
two fractional linear expressions and using \eqref{ocf:eq:det} proves
\eqref{ocf:eq:width}--\eqref{ocf:eq:distS}. The determinant also shows strict
monotonicity; direct solution for $t$ shows that the image is the entire
open interval between $r_n$ and $s_n$. Membership in every finite cylinder
is precisely indefinite continuation with the prescribed digits.
\end{proof}

\begin{corollary}[Realization in the full surreal field]\label{ocf:cor:realization}
Every admissible code has a realization in $\No$. In fact, let $L$ and $R$
be the sets of lower and upper endpoints of all the $C_n$. Then the Conway
cut $\{L\mid R\}$ belongs to $\F(\codea)$.
\end{corollary}
\begin{proof}
The sets are countable and every member of $L$ is strictly less than every
member of $R$: compare both with a point of a sufficiently late nonempty
nested interval. The standard cut property supplies a surreal strictly
between the two sets. It lies in every $C_n$.
\end{proof}

The cut in this corollary is a simplicity-based choice. The proof does not
interpret it as an order-topological limit, and later periodic fixed-point
choices will not be identified with it without a separate argument.

\section{An exact valuation description of every fiber}
\label{ocf:sec:fibers}

\subsection{Why both endpoints matter}

An upper error bound $|x-r_n|<q_n^{-2}$ does not by itself show that a
perturbation preserves the digits. One also needs a lower bound on the
distance to the \emph{other} boundary. The case of a next digit equal to
one is the delicate boundary case.

\begin{lemma}[Uniform two-step buffer]\label{ocf:lem:buffer}
For $x\in\F(\codea)$ and every $n\ge0$, the distance from $x$ to either
endpoint of $C_n$ is strictly greater than
\begin{equation}\label{ocf:eq:buffer}
 \frac1{4q_{n+2}^{\,2}}.
\end{equation}
\end{lemma}
\begin{proof}
Put $t=x_{n+1}=a+1/u$ and $u=x_{n+2}=b+1/w$, where
$a=a_{n+1}$, $b=a_{n+2}$, and $w>1$. Thus $a<t<a+1$ and $b<u<b+1$.
Since $q_nt+q_{n-1}<q_{n+1}+q_n\le2q_{n+1}$, formula
\eqref{ocf:eq:distR} gives
\[
 |x-r_n|>\frac1{2q_nq_{n+1}}\ge\frac1{2q_{n+1}^{\,2}}.
\]
For the other endpoint first suppose $a>1$. Discreteness of $\Oz$ implies
$a\ge2$, so $t-1>1$. In \eqref{ocf:eq:distS}, both denominator factors are at
most $2q_{n+1}$, and the second is strictly smaller. Consequently
$|x-s_n|>1/(4q_{n+1}^{\,2})$.

If $a=1$, then $q_{n+1}=q_n+q_{n-1}$ and $t-1=1/u$. Substitution in
\eqref{ocf:eq:distS} gives the sharper identity
\[
 |x-s_n|=\frac1{q_{n+1}(q_{n+1}u+q_n)}
 >\frac1{q_{n+1}(q_{n+2}+q_{n+1})}
 \ge\frac1{2q_{n+2}^{\,2}}.
\]
Monotonicity of the denominators now gives \eqref{ocf:eq:buffer} in all cases.
\end{proof}

\begin{lemma}[Degree and exact convergent error]\label{ocf:lem:degrees}
For every $n\ge0$,
\begin{equation}\label{ocf:eq:degrees}
 \dgr q_n=d_n=\sum_{j=1}^n\dgr a_j.
\end{equation}
For each $x\in\F(\codea)$,
\begin{equation}\label{ocf:eq:exacterror}
 v(x-r_n)=d_n+d_{n+1}.
\end{equation}
\end{lemma}
\begin{proof}
The denominator degree formula starts at $q_0=1$. In
$a_nq_{n-1}+q_{n-2}$ both summands are nonnegative and the first has degree
$\dgr a_n+d_{n-1}\ge d_{n-2}$. Leading coefficients do not cancel, so
induction proves \eqref{ocf:eq:degrees}.

In \eqref{ocf:eq:distR}, the complete quotient satisfies
$a_{n+1}<t<a_{n+1}+1$, whence $\dgr t=\dgr a_{n+1}$. The positive
expression $q_nt+q_{n-1}$ therefore has degree
$d_n+\dgr a_{n+1}=d_{n+1}$. Taking the valuation of the reciprocal in
\eqref{ocf:eq:distR} proves \eqref{ocf:eq:exacterror}.
\end{proof}

\begin{definition}[Precision ideal]\label{ocf:def:ideal}
For an admissible code let
\begin{equation}\label{ocf:eq:ideal}
 I_{\codea}=\{0\}\cup
 \{e\in\No^\times:v(e)>2d_n\text{ for all }n\ge0\}.
\end{equation}
Equivalently, $|e|\prec q_n^{-2}$ for every $n$.
\end{definition}

\begin{theorem}[Exact fiber]\label{ocf:thm:fiber}
For every admissible code $\codea$ and every $x\in\F(\codea)$,
\begin{equation}\label{ocf:eq:fiber}
 \boxed{\ \F(\codea)=x+I_{\codea}.\ }
\end{equation}
The group $I_{\codea}$ is a convex $\R$-linear additive subgroup of $\No$
and a proper ideal of $\Ofin$, contained in $\mm$.
\end{theorem}
\begin{proof}
Suppose $x,y\in\F(\codea)$. For fixed $n$ and every $k\ge0$, both points
belong to $C_{n+k}$. Thus by \eqref{ocf:eq:width} and \eqref{ocf:eq:fib},
\[
 |x-y|<q_{n+k}^{-2}\le F_{k+1}^{-2}q_n^{-2}.
\]
The ordinary Fibonacci numbers are unbounded in the ordinary integers.
For every ordinary $m$ we may choose $k$ with $F_{k+1}^2>m$. It follows
that $|x-y|\prec q_n^{-2}$. Since $n$ was arbitrary, $y-x\in I_{\codea}$.

Conversely, take $e\in I_{\codea}$. For each $n$,
$|e|\prec q_{n+2}^{-2}$, so in particular
$|e|<1/(4q_{n+2}^{\,2})$. The buffer lemma ensures $x+e\in C_n$.
This holds for all $n$, proving $x+e\in\F(\codea)$.

The valuation inequality proves additive closure, including cancellation.
Multiplication by a nonzero real does not change valuation. Multiplication
by an element of $\Ofin$ can only increase valuation. Convexity follows
either from \eqref{ocf:eq:little} or from the valuation order. Finally $d_0=0$
implies $I_{\codea}\subseteq\mm$, so $1\notin I_{\codea}$.
\end{proof}

\begin{corollary}[No unique unrestricted surreal realization]
\label{ocf:cor:neverunique}
Every admissible infinite code has a nonzero precision ideal and a
proper-class fiber in $\No$.
\end{corollary}
\begin{proof}
The set $\{2d_n:n\ge0\}$ has a strict surreal upper bound $\gamma$.
Then $\omega^{-\gamma}\in I_{\codea}$. For arbitrarily large positive
surreal $\eta>\gamma$, the distinct monomials $\omega^{-\eta}$ also lie
in the ideal. These cannot form a set: their valuations would give a
set unbounded above in $\No$, contrary to the cut property.
\end{proof}

The basic proper-class conclusion is consistent with, and already
anticipated by, the older associated-class theory. The useful refinement
here is its exact cut of allowed valuations.

\subsection{What an ordinary sequence does not converge to}

\begin{proposition}[Order-topological obstruction]\label{ocf:prop:topology}
A net indexed by a set that converges in the order topology of $\No$ is
eventually equal to its limit. In particular, the sequence of continuants
$r_n$ of an admissible infinite code does not converge in that topology.
\end{proposition}
\begin{proof}
Suppose $(z_j)_{j\in J}$ converges to $z$, where $J$ is a set. The nonzero
positive distances $|z_j-z|$ form a set $S$. If $S$ is nonempty, the cut
property gives $0<\delta<s$ for all $s\in S$. Eventual membership in
$(z-\delta,z+\delta)$ therefore means eventual equality to $z$. The
continuants are never eventually constant, since the determinant identity
makes consecutive ones distinct.
\end{proof}

This elementary class-size obstruction is not presented as a new topology
theorem. It explains why cylinder realization, coefficientwise Hahn
behavior, and order convergence must be kept separate.

\section{Realization inside a prescribed full Hahn field}
\label{ocf:sec:hahn}

We now answer a stronger question than Corollary~\ref{ocf:cor:realization}:
when the digits lie in a given full Hahn field, can the realization be kept
inside it? For an arbitrary ordered subfield of $\No$ there is no reason
for the answer to be affirmative. For $K_G$ it always is.

\begin{theorem}[Hahn realization and exact uniqueness]\label{ocf:thm:hahn}
Let $G\subseteq\No$ be a set-sized additive subgroup and let $\codea$ be
an admissible code in $\Oz_G$. Then
\begin{enumerate}[label=\textup{(\roman*)}]
\item $\F(\codea)\cap K_G\ne\varnothing$;
\item for $x\in\F(\codea)\cap K_G$,
\[
 \F(\codea)\cap K_G=x+(I_{\codea}\cap K_G);
\]
\item this fiber is a singleton if and only if $(2d_n)$ is cofinal in
$G_{\ge0}$, in the sense that every $g\in G_{\ge0}$ satisfies
$g\le2d_n$ for some $n$.
\end{enumerate}
No divisibility assumption on $G$ is needed.
\end{theorem}
\begin{proof}
Only existence requires a new argument. All continuants $r_n$ are already
in $K_G$. Split the proof according to the degree sequence.

\emph{Case 1: $d_n$ is eventually constant.} Each increment
$d_n-d_{n-1}=\dgr a_n$ is nonnegative. Thus all sufficiently late digits
have degree zero and are positive ordinary integers. The corresponding
ordinary real infinite continued fraction has an irrational real value
$\rho>1$ \cite{Krishnan}. Apply the finite prefix map $\Phi_{N-1}$ to
$\rho$, where the ordinary tail starts at $a_N$ and $N\ge1$. This produces
a point of $K_G$ with the required code. The denominator is positive, and
all tail remainders are nonzero by the ordinary irrationality theorem.

\emph{Case 2: $d_n$ is not eventually constant.} Choose any
$x\in\F(\codea)$ in $\No$. Delete from its normal form the terms whose
growth exponents are outside $G$, and call the resulting subseries
$\pi_G(x)$. A subseries of a reverse well-ordered support is allowed, so
$\pi_G(x)\in K_G$.

Fix $n$. There is $k$ with $d_k>d_n$. By \eqref{ocf:eq:exacterror},
\[
 v(x-r_k)=d_k+d_{k+1}>2d_n.
\]
Since $r_k$ is supported in $G$, every term of $x$ outside $G$ is also a
term of $x-r_k$. Every deleted term therefore has valuation greater than
$2d_n$. This is true for every $n$, and if the deleted subseries is nonzero
its leading term is one of those terms. Hence
\[
 x-\pi_G(x)\in I_{\codea}.
\]
The exact fiber theorem implies $\pi_G(x)\in\F(\codea)$, proving existence.

Assertion (ii) is the intersection of \eqref{ocf:eq:fiber} with $K_G$.
For (iii), a nonzero element of $I_{\codea}\cap K_G$ has a valuation
$g\in G$ strictly above every $2d_n$. Conversely every such $g$ is realized
by $\omega^{-g}\in K_G$. This is exactly the failure of cofinality.
\end{proof}

\begin{remark}[The projection is not a homomorphism]\label{ocf:rem:projection}
The map $\pi_G$ in the proof is an additive support projection, not a field
map. For $G=\Z$, one has
$\pi_G(\omega^{1/2})=0$ but $\pi_G((\omega^{1/2})^2)=\omega$.
The argument uses only the support of the difference from an actual
continuant; it never projects a product or a reciprocal.
\end{remark}

\begin{remark}[Why eventual constancy was separated]\label{ocf:rem:twocases}
It would be invalid to demand a $k$ with $d_k>d_n$ in the first case.
Ordinary digits improve real coefficient precision without changing
valuation. The real continued-fraction theorem is exactly the missing
input in that case. The two-case proof also avoids an unjustified
summation of the differences $r_{n+1}-r_n$.
\end{remark}

\begin{corollary}[Countable cofinality threshold]\label{ocf:cor:cofinality}
For a nonzero set-sized additive subgroup $G$ of $\No$, some admissible
infinite code in $\Oz_G$ has a unique realization in $K_G$ if and only if
$G_{\ge0}$ has a countable cofinal subset.
\end{corollary}
\begin{proof}
A uniquely realizing code supplies the countable cofinal sequence $(2d_n)$.
Conversely choose a nondecreasing cofinal sequence $g_n\in G_{\ge0}$,
with $g_0=0$. Set
\[
 a_0=0,\qquad a_n=\omega^{g_n-g_{n-1}}\quad(n\ge1).
\]
Zero increments give the digit $1$; positive increments give a positive
infinite omnific integer. The degree sequence is exactly $g_n$.
Theorem~\ref{ocf:thm:hahn} gives a unique realization.
\end{proof}

The cofinality statements using $(d_n)$ and $(2d_n)$ are equivalent. One
direction is immediate; for the other, apply cofinality of $(2d_n)$ to
$2g$ to obtain $d_n\ge g$. This does not divide an element of $G$ by two.

\begin{corollary}[Convergence in the unique Hahn case]\label{ocf:cor:convergence}
If $G\ne\{0\}$ and the conditions for uniqueness in
Theorem~\ref{ocf:thm:hahn} hold, the continuants converge in the order topology
of $K_G$ to that unique point.
\end{corollary}
\begin{proof}
For any positive $b\in K_G$, choose $n_0$ with
$d_{n_0}+d_{n_0+1}>v(b)$. Such a strict bound exists by cofinality and a
positive element of $G$. For $n\ge n_0$, \eqref{ocf:eq:exacterror} gives
$|x-r_n|\prec b$, hence $|x-r_n|<b$.
\end{proof}

When $G=\{0\}$, the field is $\R$ and ordinary real convergence applies
instead. When $G\ne\{0\}$ but all late digits are ordinary, ordinary
real coefficient convergence does \emph{not} imply order convergence in
$K_G$: nonzero real errors exceed every positive infinitesimal.

\section{The algebra of precision ideals}
\label{ocf:sec:ideals}

\subsection{Realizing an arbitrary countably specified precision cut}

\begin{proposition}[Prescribed degree cut]\label{ocf:prop:prescribed}
Let $(\gamma_n)_{n\ge0}$ be a nondecreasing sequence of nonnegative
surreals with $\gamma_0=0$. There is an admissible code with
\[
 I_{\codea}=\{0\}\cup\{e:v(e)>\gamma_n\text{ for every }n\}.
\]
One choice is
\begin{equation}\label{ocf:eq:prescribed}
 a_0=0,\qquad a_n=\omega^{(\gamma_n-\gamma_{n-1})/2}\quad(n\ge1).
\end{equation}
\end{proposition}
\begin{proof}
Every exponent in \eqref{ocf:eq:prescribed} is nonnegative, so each digit is at
least one and belongs to $\Oz$. Formula \eqref{ocf:eq:degrees} gives
$2d_n=\gamma_n$, which proves the assertion.
\end{proof}

Within a prescribed $G$, this construction requires the displayed half
increments to belong to $G$. This is automatic for divisible $G$ when all
$\gamma_n\in G$. Theorem~\ref{ocf:thm:hahn} itself has no such divisibility
requirement.

\begin{theorem}[Prime precision criterion]\label{ocf:thm:prime}
The ideal $I_{\codea}$ is prime in $\Ofin$ if and only if condition
\eqref{ocf:cond:P} holds. The same equivalence holds for
$I_{\codea}\cap K_G$ in $\Ofin\cap K_G$ when the digits lie in $K_G$.
If \eqref{ocf:cond:P} fails, there is an explicit element whose image in the
quotient is nonzero but has square zero.
\end{theorem}
\begin{proof}
The nonnegative valuations of elements outside $I_{\codea}$ form the lower
set
\[
 D=\bigcup_{n\ge0}[0,2d_n].
\]
Primality is equivalent to closure of $D$ under addition. If
\eqref{ocf:cond:P} holds and $g,h\in D$, choose $n$ bounding both by $2d_n$
and then $m$ with $d_m\ge2d_n$. It follows that
$g+h\le4d_n\le2d_m$, so $g+h\in D$.

Conversely, additive closure applied to $2d_n+2d_n$ gives an $m$ with
$4d_n\le2d_m$, which is \eqref{ocf:cond:P}.
If this condition fails, choose $n$ such that $d_m<2d_n$ for every $m$.
Then $d_n>0$ and
\[
 e=\omega^{-2d_n}\notin I_{\codea},\qquad
 e^2=\omega^{-4d_n}\in I_{\codea}.
\]
When all digits lie in $K_G$, these monomials lie there too; thus the same
necessity proof works in the restricted ring. The sufficiency proof is
unchanged.
\end{proof}

Primality for the full class is an elementwise assertion. Quotient rings
are taken inside set-sized Hahn workspaces containing the digits and the
elements under discussion. Every displayed witness lies in such a
workspace. No object consisting of proper-class cosets is assumed.

\begin{corollary}[Radical and nilpotency index]\label{ocf:cor:radical}
For $0\ne e\in\Ofin$, its class modulo $I_{\codea}$ in a set-sized
workspace as above is nilpotent exactly when
\begin{equation}\label{ocf:eq:nilpotency}
 \text{there is }k\ge1\text{ such that }k\,v(e)>2d_n
 \text{ for every }n.
\end{equation}
The least such $k$ is its nilpotency index, with index one meaning the zero
class. The radical $\sqrt{I_{\codea}}$ is prime.
\end{corollary}
\begin{proof}
The criterion and index are the equality $v(e^k)=kv(e)$ and the definition
of the ideal. To prove primality, suppose $(uw)^k\in I_{\codea}$.
Write $g=v(u)$, $h=v(w)$ and assume $g\le h$. Then
$2kh\ge k(g+h)>2d_n$ for every $n$, so $w^{2k}\in I_{\codea}$ and
$w\in\sqrt{I_{\codea}}$. Zero factors are immediate. This is the familiar
valuation-ring argument, included to identify the precise cut involved.
\end{proof}

Condition \eqref{ocf:cond:P} is thus more than a growth convenience. It says
that the undetected infinitesimals define a reduced integral-domain
quotient, rather than an error calculus with nonzero nilpotent classes.
It will also be one of the two exact conditions for algebraic rigidity.

\section{Algebraic rigidity over the digit field}
\label{ocf:sec:algebraic}

\subsection{The coefficient field must be specified}

Define
\[
 K_{\codea}=\Q(a_0,a_1,a_2,\ldots),\qquad
 A_{\codea}=\{x\in\F(\codea):x\text{ is algebraic over }K_{\codea}\}.
\]
The field $K_{\codea}$ is countable: every element is a rational expression
in finitely many members of a countable list. The digits themselves may
have enormous or uncountable supports, but this does not change the
cardinality of the field \emph{generated by them}. Consequently the set of
all elements of $\No$ algebraic over $K_{\codea}$ is countable.

These are statements relative to the specified digit field, not to all
omnific coefficients or to the whole ambient surreal field. For an
arbitrary code we do not know that $A_{\codea}$ is nonempty.

\begin{lemma}[Algebraic errors cannot be newly invisible]
\label{ocf:lem:algerror}
Let $E\subseteq\No$ be a set-sized subfield. If a nonzero
$e\in I_{\codea}$ is algebraic over $E$, then
$I_{\codea}\cap E$ contains a nonzero element. Hence
\[
 I_{\codea}\cap E=\{0\}
 \quad\Longleftrightarrow\quad
 I_{\codea}\cap\alg_{\No}(E)=\{0\},
\]
where $\alg_{\No}(E)$ is the field of elements of $\No$ algebraic over $E$.
\end{lemma}
\begin{proof}
Put $g=v(e)>0$ and take a nonzero polynomial relation
$\sum_j c_je^j=0$ over $E$. Among the valuations $v(c_j)+jg$ of nonzero
terms, the minimum occurs at least twice. Otherwise the uniquely leading
term could not cancel. For two minimizing indices $i<j$,
\[
 (j-i)g=v(c_i)-v(c_j)=v(c_i/c_j).
\]
The ratio $c_i/c_j$ is a nonzero element of $E$. Since $g>2d_n$ for all
$n$, its valuation $(j-i)g$ has the same strict lower bounds. Thus the
ratio belongs to $I_{\codea}$. The reverse implication in the equivalence
is immediate from $E\subseteq\alg_{\No}(E)$.
\end{proof}

\begin{lemma}[What the digit field can see]\label{ocf:lem:digitfield}
With $D_0=\dgr a_0$ for $a_0\ne0$, and $D_0=0$ otherwise,
\begin{equation}\label{ocf:eq:fieldvisibility}
 I_{\codea}\cap K_{\codea}=\{0\}
 \quad\Longleftrightarrow\quad
 \eqref{ocf:cond:P}\text{ and }\eqref{ocf:cond:V}.
\end{equation}
If either condition fails, a nonzero element of the intersection is
explicitly available from the digits and continuants.
\end{lemma}
\begin{proof}
If \eqref{ocf:cond:V} fails, then $a_0\ne0$ and
$v(a_0^{-1})=D_0>2d_n$ for every $n$. Thus $a_0^{-1}\in I_{\codea}\cap
K_{\codea}$.
If \eqref{ocf:cond:P} fails, choose $n$ with $d_m<2d_n$ for every $m$.
The nonzero element $q_n^{-4}\in K_{\codea}$ has valuation $4d_n>2d_m$
for every $m$, so it belongs to the intersection.

For the converse assume both conditions. Let $z\in K_{\codea}^\times$.
Write $z=P/Q$, where $P,Q$ are nonzero evaluations of integer-coefficient
polynomials in finitely many digits. They belong to $\Oz$. Choose $N$
large enough to include those digit indices and to satisfy $D_0\le2d_N$.
Every degree of an appearing digit is then at most $2d_N$.
If $M$ bounds the total degrees of both polynomials, every nonzero
monomial in them has growth degree at most $2Md_N$. Finite addition,
even with cancellation, cannot increase that bound. Therefore
\[
 0\le\dgr P,\dgr Q\le2Md_N,
 \qquad |v(z)|=|\dgr Q-\dgr P|\le2Md_N.
\]
If $d_N=0$, this gives $v(z)=0$, already excluding $z\in I_{\codea}$.
Otherwise iterating \eqref{ocf:cond:P} produces an $m$ with
$d_m\ge M d_N$. It follows that $v(z)\le2d_m$, again excluding membership.
This proves \eqref{ocf:eq:fieldvisibility}.
\end{proof}

\begin{theorem}[Exact algebraic-fiber trichotomy]\label{ocf:thm:algebraic}
For every admissible code the following statements hold.
\begin{enumerate}[label=\textup{(\roman*)}]
\item If \eqref{ocf:cond:P} and \eqref{ocf:cond:V} both hold, then
$|A_{\codea}|\le1$.
\item If either condition fails, then $A_{\codea}$ is either empty or
countably infinite.
\item In the nonempty case of (ii), for any $\xi\in A_{\codea}$ there is
$0\ne\delta\in I_{\codea}\cap K_{\codea}$ such that
\begin{equation}\label{ocf:eq:algtranslations}
 \xi+q\delta\in A_{\codea},\qquad
 K_{\codea}(\xi+q\delta)=K_{\codea}(\xi)
 \quad(q\in\Q).
\end{equation}
\end{enumerate}
Thus whenever an algebraic realization exists, it is unique if and only
if both displayed conditions hold.
\end{theorem}
\begin{proof}
If $\xi,\eta$ are algebraic realizations, their difference is algebraic
over $K_{\codea}$ and belongs to $I_{\codea}$ by
Theorem~\ref{ocf:thm:fiber}. Lemmas~\ref{ocf:lem:algerror} and
\ref{ocf:lem:digitfield} force that difference to vanish when both conditions
hold.

When either fails, Lemma~\ref{ocf:lem:digitfield} supplies a nonzero $\delta$
in the digit field and precision ideal. Given one algebraic $\xi$, the
exact fiber theorem and $\R$-linearity of the ideal give all the points
in \eqref{ocf:eq:algtranslations}. Translation by an element of the base field
does not change the generated extension. The points are distinct as $q$
varies. Countability of the entire algebraic closure inside $\No$ gives
the upper bound and hence the exact countable cardinality.
\end{proof}

\subsection{Two different sources of failure}

Condition \eqref{ocf:cond:P} is intrinsic to the precision ideal: the observed
denominator degrees must eventually absorb their own doubles. Condition
\eqref{ocf:cond:V} is different. The first digit contributes a parameter to
the coefficient field but contributes nothing to any denominator degree.
It may therefore introduce a scale invisible to all subsequent digits.

This distinction explains why even a long periodic tail can fail to give
algebraic uniqueness. It also explains why the full-class nonuniqueness
of Corollary~\ref{ocf:cor:neverunique} is compatible with uniqueness among
algebraic elements over a countable coefficient field.

\section{Periodic itineraries and periodic complete quotients}
\label{ocf:sec:periodic}

\subsection{A distinguished algebraic realization always exists}

Assume an eventually periodic code is displayed as
\begin{equation}\label{ocf:eq:periodiccode}
 \codea=(a_0,\ldots,a_{r-1},\overline{b_1,\ldots,b_\ell}),
 \qquad r\ge1,\quad\ell\ge1,
\end{equation}
where the $b_j$ are positive omnific integers. The finite digit field is
\[
 K=\Q(a_0,\ldots,a_{r-1},b_1,\ldots,b_\ell)=K_{\codea}.
\]
Let
\begin{equation}\label{ocf:eq:periodmatrix}
 M=\prod_{j=1}^{\ell}\begin{pmatrix}b_j&1\\1&0\end{pmatrix}
   =\begin{pmatrix}A&B\\C&D\end{pmatrix},
 \qquad T(t)=\frac{At+B}{Ct+D}.
\end{equation}
Here $B,C>0$, $A,D\ge0$, and $\det M=(-1)^\ell$.

\begin{theorem}[Period-fixed representative]\label{ocf:thm:periodroot}
There is exactly one $t>1$ satisfying $T(t)=t$. It is the positive root of
\begin{equation}\label{ocf:eq:periodpoly}
 Ct^2+(D-A)t-B=0,
\end{equation}
namely
\begin{equation}\label{ocf:eq:periodroot}
 t=\frac{A-D+\sqrt{(D-A)^2+4BC}}{2C}.
\end{equation}
Its complete quotients repeat after the block. The point
$\xi=\Phi_{r-1}(t)$ belongs to $\F(\codea)$ and satisfies $[K(\xi):K]\le2$.
It is the unique point in that fiber whose complete quotient at index $r$
is fixed by this period transformation.
\end{theorem}
\begin{proof}
The discriminant is positive and $\No$ is real closed. The two roots have
product $-B/C<0$, so exactly one is positive. Moreover $T(1)>1$: a finite
composition of the maps $u\mapsto b_j+1/u$, applied at $u=1$, has first
value strictly greater than $b_1\ge1$. Thus the polynomial in
\eqref{ocf:eq:periodpoly} is negative at $1$, and its positive root exceeds
$1$.

For this root, the identity $t=T(t)$ presents one full block of digits
with a final tail equal to $t>1$. All intermediate tails are greater than
one, so the floor algorithm produces the block and returns to the same
complete quotient. Induction repeats this indefinitely. Applying the
finite prefix gives $\xi$ with the code \eqref{ocf:eq:periodiccode}. The
quadratic equation and the fractional linear prefix place $\xi$ in an
extension of degree at most two. Finally, the prefix map is injective,
and \eqref{ocf:eq:periodpoly} has only one root greater than one.
\end{proof}

This is the familiar fixed-point mechanism behind periodic continued
fractions, extended by ordered-field algebra. It is not the assertion
that \emph{every} point with periodic digits satisfies
\eqref{ocf:eq:periodpoly}. A periodic itinerary does not force its complete
quotients to recur as equal numbers.

\subsection{The exact algebraic dichotomy}

Define
\begin{equation}\label{ocf:eq:periodscales}
\begin{split}
 E&=\sum_{j=1}^{\ell}\dgr b_j,\\
 D_*&=\max\bigl(\{0,D_0\}\cup
 \{\dgr a_j:1\le j<r\}\cup\{\dgr b_j:1\le j\le\ell\}\bigr).
\end{split}
\end{equation}
Thus $E$ is the degree gained per full period. It is not the degree of a
matrix entry chosen without regard to the recurrence.

\begin{theorem}[One or countably many algebraic periodic realizations]
\label{ocf:thm:periodiccriterion}
For the code \eqref{ocf:eq:periodiccode}, exactly one point of $\F(\codea)$ is
algebraic over $K$ if and only if one of the following holds:
\begin{equation}\label{ocf:eq:periodcriterion}
 \boxed{\quad E=D_*=0\quad\text{or}\quad
 E>0\ \text{and}\ D_*\le N E
 \text{ for some ordinary }N\ge1.\quad}
\end{equation}
In that case the point is $\xi$ from Theorem~\ref{ocf:thm:periodroot}.
Otherwise there are countably infinitely many algebraic realizations.
In particular there are infinitely many of degree at most two, all lying
in the single extension $K(\xi)$.
\end{theorem}
\begin{proof}
Put $S=\sum_{j=1}^{r-1}\dgr a_j$ and
$E_t=\sum_{j=1}^{t}\dgr b_j$, with $E_0=0$. For $k\ge0$ and
$0\le t<\ell$,
\begin{equation}\label{ocf:eq:perioddegree}
 d_{r-1+k\ell+t}=S+kE+E_t.
\end{equation}
If $E=D_*=0$, both conditions \eqref{ocf:cond:P} and \eqref{ocf:cond:V} hold.
If $E>0$ and $D_*\le NE$, all prefix degrees and all partial-period sums
are bounded by ordinary multiples of $E$. The unbounded ordinary
multiples $kE$ in \eqref{ocf:eq:perioddegree} then absorb both the double of
any specified $d_n$ and $D_0$. Again both conditions hold. Theorem
\ref{ocf:thm:algebraic} gives at most one algebraic point and
Theorem~\ref{ocf:thm:periodroot} supplies it.

Suppose \eqref{ocf:eq:periodcriterion} fails. Then $D_*>0$, and either $E=0$
or $D_*>NE$ for every ordinary $N$. Choose a displayed nonzero digit $b$
with degree $D_*$, allowing $b=a_0$ if necessary, and put
\[
 \delta=b^{-(2r+2)}\in K^\times.
\]
Its valuation is $(2r+2)D_*$. The prefix bound
$S\le(r-1)D_*$ and \eqref{ocf:eq:perioddegree} show that this valuation exceeds
$2d_n$ for every $n$. For $E>0$, use
$2(k+1)E<D_*$ for each ordinary $k$; for $E=0$ the assertion is immediate.
Thus $\delta\in I_{\codea}\cap K$. The points $\xi+q\delta$, $q\in\Q$,
are distinct algebraic realizations and generate exactly $K(\xi)$.
Countability gives the asserted total cardinality of algebraic points.
\end{proof}

\begin{corollary}[Pure periodicity has one algebraic representative]
\label{ocf:cor:pureperiod}
For a purely periodic admissible code with positive first digit, exactly
one realization is algebraic over the field generated by its period
block. All other realizations are transcendental over that field.
\end{corollary}
\begin{proof}
Display the code with a one-digit prefix and a cyclically rotated block.
Every displayed digit occurs in the period, so $D_*\le E$ when $E>0$.
If $E=0$, all digits are ordinary and $D_*=0$. Apply
Theorem~\ref{ocf:thm:periodiccriterion}.
\end{proof}

\subsection{A stricter cofinality threshold for periodic codes}

A positive element $e$ of an ordered abelian group $G$ is an \emph{order
unit} if every $g\in G_{\ge0}$ is bounded above by an ordinary multiple
of $e$.

\begin{theorem}[Periodic uniqueness detects an order unit]
\label{ocf:thm:orderunit}
Let $G\ne\{0\}$ and let \eqref{ocf:eq:periodiccode} be a code in $\Oz_G$.
Its realization in $K_G$ is unique if and only if its period degree $E$
is a positive order unit of $G$. Consequently:
\begin{align*}
 &\text{some infinite code is unique in }K_G
   &&\Longleftrightarrow\quad G\text{ has countable cofinality},\\
 &\text{some eventually periodic code is unique in }K_G
   &&\Longleftrightarrow\quad G\text{ has a positive order unit}.
\end{align*}
\end{theorem}
\begin{proof}
If $E$ is an order unit, the terms $S+kE$ in \eqref{ocf:eq:perioddegree} are
cofinal in $G_{\ge0}$. Theorem~\ref{ocf:thm:hahn} gives uniqueness.

Conversely, uniqueness makes $(d_n)$ cofinal. It cannot happen with
$E=0$, because then the degrees are eventually bounded by $S$. Given
$g\in G_{\ge0}$, choose a sufficiently late $n=r-1+k\ell+t$ with
$d_n\ge S+g$. Formula \eqref{ocf:eq:perioddegree} implies
$g\le kE+E_t\le(k+1)E$. Thus $E$ is an order unit.

The first existence equivalence is Corollary~\ref{ocf:cor:cofinality}. For the
second, an order unit $e$ gives the constant code
$(\omega^e,\omega^e,\ldots)$, whose degree sequence is $d_n=ne$.
Necessity follows from the preceding paragraph.
\end{proof}

This separates two notions often conflated by rank-one intuition. A value
group can have countable cofinality without any order unit. In such a
field some infinite codes determine unique points, but no eventually
periodic code does so.

\wtag{} The same order-unit threshold decides, in the expanding-polynomial
report, whether an itinerary word determines one point
(\lbl{epd:thm:fibers}); Section~\ref{ocf:sec:collection} compares the two
fibre ideals.

\section{Worked examples and boundary cases}
\label{ocf:sec:examples}

\subsection{The golden ratio and its whole infinitesimal neighborhood}

For $\codea=(1,1,1,\ldots)$ all $d_n=0$. Hence
\begin{equation}\label{ocf:eq:golden}
 \F(\codea)=\varphi+\mm,
 \qquad \varphi=\frac{1+\sqrt5}{2}.
\end{equation}
The algebraic representative over $K=\Q$ is uniquely $\varphi$. For
example, $\varphi+\omega^{-1}$ has exactly the same infinite digits but is
not algebraic over $\Q$. Indeed it is a nonreal surreal, and no nonreal
surreal is algebraic over the real-closed ordered subfield $\R$.
The example illustrates the older nonuniqueness phenomenon; the new
classification explains its algebraic rigidity at the same time.

In $K_{\{0\}}=\R$ the fiber is a singleton. In every $K_G$ with
$G\ne\{0\}$ it is not, since $K_G$ has a positive infinitesimal. Ordinary
real continued-fraction uniqueness is therefore exactly consistent with
\eqref{ocf:eq:golden}.

\subsection{An initial scale that a periodic tail never sees}

For $\codea=(\omega,1,1,1,\ldots)$ again $I_{\codea}=\mm$, but now
$K=\Q(\omega)$. A period-fixed realization is
\[
 \xi=\omega+\frac1\varphi=\omega+\varphi-1.
\]
Every $\xi+q/\omega$, $q\in\Q$, has the same code and is algebraic of
degree two over $K$. Here \eqref{ocf:cond:P} holds, whereas
\eqref{ocf:cond:V} fails: $D_0=1$ but every $d_n=0$.
Thus even primeness of the precision ideal does not imply algebraic
uniqueness over the digit field.

\subsection{One infinite later digit and a nonreduced error quotient}

Take $\codea=(0,\omega,1,1,\ldots)$. Then
$d_0=0$ and $d_n=1$ for $n\ge1$, so
\[
 I_{\codea}=\{0\}\cup\{e:v(e)>2\}.
\]
The element $\omega^{-2}$ is outside the ideal but its square is inside.
The quotient therefore has a nonzero square-zero class. Here
\eqref{ocf:cond:V} holds and \eqref{ocf:cond:P} fails. A distinguished point is
$\xi=1/(\omega+1/\varphi)$, and $\xi+q\omega^{-3}$ gives infinitely many
algebraic realizations over $\Q(\omega)$.

\subsection{A constant infinite period}

For $\codea=(\omega,\omega,\omega,\ldots)$ the positive fixed point is
\begin{equation}\label{ocf:eq:rho}
\begin{split}
 \rho&=\frac{\omega+\sqrt{\omega^2+4}}2\\
 &=\omega+\omega^{-1}-\omega^{-3}+2\omega^{-5}
  -5\omega^{-7}+14\omega^{-9}-\cdots.
\end{split}
\end{equation}
The binomial series at the infinitesimal $\omega^{-2}$ gives the displayed
Hahn expansion. More precisely its coefficient of $\omega^{-2k-1}$ is
$(-1)^kC_k$, where $C_k=\frac1{k+1}\binom{2k}{k}$.
Here $d_n=n$, and
\[
 I_{\codea}=\{0\}\cup\{e:v(e)>2n\text{ for every ordinary }n\}.
\]
There is a unique realization in $\R((\omega^\Q))$ and in
$\R((\omega^\R))$, but not in $\No$: the perturbation
$\omega^{-\omega}$ is invisible to the entire code. Nevertheless $\rho$
is the unique realization algebraic over $\Q(\omega)$.

The code is periodic, its full-class fiber is large, its rank-one Hahn
fiber is a singleton, and its digit-field algebraic fiber is a singleton.
These four statements concern different ambient categories and do not
contradict one another.

\subsection{A bounded increasing cut, with a genuinely different boundary}

Let $a_0=0$ and $a_n=\omega^{2^{-n}}$ for $n\ge1$. Then
\[
 d_n=1-2^{-n},\qquad
 I_{\codea}=\{0\}\cup\{e:v(e)>2-2^{1-n}\text{ for every }n\}.
\]
In the rank-one field $K_\R$, this becomes
\begin{equation}\label{ocf:eq:rankoneboundary}
 I_{\codea}\cap K_\R=\{0\}\cup\{e\in K_\R:v(e)\ge2\}.
\end{equation}
The endpoint is included. In the full surreal value group one must retain
the sequence of inequalities: the value $2-\omega^{-1}$ satisfies them
all but is less than $2$. Replacing a real supremum by a nonexistent
surreal supremum would give the wrong ideal.

Theorem~\ref{ocf:thm:hahn} ensures a realization in $K_\R$ even though the
cumulative degrees are bounded and the fiber is not a singleton.
This example exercises the non-eventually-constant support-projection
case, not the ordinary-tail case.

\subsection{Countable cofinality without an order unit}

Let
\[
 G=\bigoplus_{j\ge1}\Q\,\omega^j\ \subseteq(\No,+),
\]
where the direct sum means finite rational linear combinations. This
ordered group has countable cofinality, witnessed by the sequence
$\omega^j$, but has no order unit: a finite linear combination is dominated
by a sufficiently higher power of $\omega$.

The code $a_0=0$, $a_n=\omega^{\omega^n}$ has cumulative degrees
$d_n=\sum_{j=1}^n\omega^j$, cofinal in $G$. It therefore has a unique
realization in $K_G$. On the other hand,
Theorem~\ref{ocf:thm:orderunit} excludes uniqueness for every eventually
periodic code in this field. This is an explicit infinite-rank separation
between the two uniqueness thresholds.

\section{Arbitrarily large algebraically independent families in a fiber}
\label{ocf:sec:independence}

The size of a full surreal fiber is not merely a counting defect. Its
undetected directions can carry arbitrarily much algebraic information.
The following statement is deliberately set-indexed, so it does not
presuppose a proper-class transcendence basis.

\begin{theorem}[Hidden algebraic independence]\label{ocf:thm:independence}
Let $\codea$ be any admissible infinite code, let $E\subseteq\No$ be a
set-sized subfield, and let $\kappa$ be any set cardinal. Then
$\F(\codea)$ contains a family of $\kappa$ elements algebraically
independent over $E$.
\end{theorem}
\begin{proof}
Choose $\xi\in\F(\codea)$. Let $H$ be the rational span of all growth
exponents occurring in the normal forms of $\xi$ and all members of $E$.
This is a set-sized subgroup of $\No$, and
$E\cup\{\xi\}\subseteq\R((\omega^H))$.
Choose an ordinal $\beta$ so large that
\begin{equation}\label{ocf:eq:betabound}
 |h|<\omega^\beta\quad(h\in H),\qquad
 2d_n<\omega^\beta\quad(n\ge0).
\end{equation}
Such a $\beta$ exists because the quantities to be bounded form a set and
ordinals are cofinal in $\No$.

For $\alpha<\kappa$ put
\[
 \theta_\alpha=\omega^{\beta+\alpha+1},\qquad
 \delta_\alpha=\omega^{-\theta_\alpha}.
\]
Each $\delta_\alpha$ belongs to $I_{\codea}$ by \eqref{ocf:eq:betabound}.
Moreover the family $(\theta_\alpha)$ is rationally independent modulo
$H$: a nonzero finite rational linear combination has a largest growth
exponent at least $\beta+1$, and hence absolute magnitude greater than
$\omega^\beta$. It cannot lie in $H$.

Consider a nonzero polynomial in finitely many $\delta_\alpha$ over
$\R((\omega^H))$. A monomial with multi-index $\nu$ shifts the support of
its nonzero coefficient into the coset
\[
 H-\sum_\alpha\nu_\alpha\theta_\alpha.
\]
Distinct multi-indices give distinct cosets by the independence just
proved. The nonzero polynomial summands consequently have disjoint
supports and cannot cancel. Thus the $\delta_\alpha$ are algebraically
independent over $\R((\omega^H))$.

Finally set $x_\alpha=\xi+\delta_\alpha$. All these points lie in the
fiber. Translation of indeterminates by $\xi$, which belongs to the
coefficient Hahn field, preserves nonzero polynomials. The $x_\alpha$
are therefore algebraically independent over $E$ as claimed.
\end{proof}

\begin{remark}[What this does and does not imply]\label{ocf:rem:independence}
The theorem holds for each prescribed set field $E$ and each prescribed
set cardinal separately. It does not choose one family independent over
\emph{all} set subfields simultaneously. Such a reading would be false,
since a field could contain a chosen member. Nor does it claim a
first-order definable or algorithmically computable independent family.
\end{remark}

\wtag{} The support-coset argument in the proof is the standard mechanism
of the independent-copies and bounded-support reports
(\lbl{isc:lem:slice}, \lbl{isc:thm:sliceddisjoint}, \lbl{bst:eq:projection});
it is credited to them, and the source's direct proof is kept
(Section~\ref{ocf:sec:collection}).

\section{Surcomplex fibers and changes of coordinate}
\label{ocf:sec:complex}

\subsection{Coordinatewise coding}

There is no order-compatible floor on $\SC$. We therefore specify the
construction rather than silently importing real continued fractions.
For $z=x+iy$, record the omnific floor code of $x$ and the omnific floor
code of $y$, whenever both algorithms are nonterminating.

\begin{theorem}[Exact surcomplex coordinate fiber]\label{ocf:thm:complex}
If $x\in\F(\codea)$ and $y\in\F(\codeb)$, then the points with the same
pair of coordinate codes form exactly
\begin{equation}\label{ocf:eq:complexfiber}
 (x+iy)+(I_{\codea}+iI_{\codeb}).
\end{equation}
If $I_{\codea}=I_{\codeb}=I$, then $I+iI$ is an ideal of
$\Ofin[i]=\Ofin+i\Ofin$. With
\[
 v_\C(u+iw)=\min\{v(u),v(w)\},
\]
it is the corresponding surcomplex valuation ideal. If the two ideals
are different, their sum in \eqref{ocf:eq:complexfiber} is not invariant under
multiplication by $i$.
\end{theorem}
\begin{proof}
The exact real fiber theorem applied to both coordinates proves
\eqref{ocf:eq:complexfiber}. If the ideals agree, then
\[
 (a+ib)(e+if)=(ae-bf)+i(af+be)\in I+iI
 \quad(a,b\in\Ofin,\ e,f\in I).
\]
The complex normal form combines the real and imaginary coefficients;
a nonzero real coefficient and a nonzero imaginary coefficient cannot
cancel each other. This proves the displayed valuation formula and the
valuation-ideal assertion.
If multiplication by $i$ preserved $I_{\codea}+iI_{\codeb}$, applying it to
real and imaginary elements would force both inclusions
$I_{\codea}\subseteq I_{\codeb}$ and $I_{\codeb}\subseteq I_{\codea}$.
\end{proof}

Equal cuts therefore give a cut-defined intersection of valuation balls,
whereas unequal cuts give an anisotropic coordinate neighborhood. This is a coordinatewise theory,
not a theorem about nearest-Gaussian-integer continued fractions or a
rotation-invariant coding of all surcomplex points.

\begin{corollary}[Surcomplex hidden independence]\label{ocf:cor:complexindependence}
Given a pair of admissible infinite coordinate codes, a set-sized
subfield $E\subseteq\SC$, and any set cardinal $\kappa$, their common
coordinate fiber contains $\kappa$ points algebraically independent over
$E$.
\end{corollary}
\begin{proof}
Fix $z=x+iy$ in the fiber. Gather the growth supports of the real and
imaginary parts of $z$ and of all members of $E$ into a rational span $H$.
Use the construction of Theorem~\ref{ocf:thm:independence}, bounding the degree
cuts of both coordinate codes, and work over $\C((\omega^H))$ instead of
$\R((\omega^H))$. The real perturbations $\delta_\alpha$ are invisible in
the real coordinate, and the imaginary coordinate is unchanged. Distinct
support cosets again prevent cancellation, including complex coefficients.
\end{proof}

\subsection{Exact transport under a fractional linear map}

\begin{proposition}[Precision transport]\label{ocf:prop:mobius}
Let $f(z)=(az+b)/(cz+d)$ over $\SC$, with $ad-bc\ne0$ and $cz+d\ne0$.
For a nonzero perturbation $e$ satisfying
\begin{equation}\label{ocf:eq:smalltransport}
 v_\C(ce)>v_\C(cz+d),
\end{equation}
one has $c(z+e)+d\ne0$ and
\begin{align}
 f(z+e)-f(z)
   &=\frac{(ad-bc)e}{(cz+d)(c(z+e)+d)},\label{ocf:eq:mobiusdiff}\\
 v_\C(f(z+e)-f(z))
   &=v_\C(ad-bc)+v_\C(e)-2v_\C(cz+d).\label{ocf:eq:mobiusscale}
\end{align}
The hypothesis is automatic when $c=0$, with $v_\C(0)=+\infty$.
\end{proposition}
\begin{proof}
The strict inequality makes $c(z+e)+d$ have the same leading valuation as
$cz+d$. Subtracting the two fractions proves \eqref{ocf:eq:mobiusdiff}; taking
valuations proves \eqref{ocf:eq:mobiusscale}.
\end{proof}

Thus within a sufficiently small fiber neighborhood, a projective change
of chart translates the valuation threshold by an exact amount. Without
\eqref{ocf:eq:smalltransport}, a pole or leading cancellation can invalidate
that linear threshold law. No claim is made that an arbitrary fractional
linear map preserves the digit codes themselves.

\section{Proof dependencies and reproducible checks}
\label{ocf:sec:verification}

The mathematical dependency chain is short. Normal forms and the omnific
floor give the finite matrix identities and cylinders. The two-step buffer
and Fibonacci growth give the exact fiber. The exact error valuation then
gives full Hahn-field realization by support projection. Independently,
the valuation cut gives the prime criterion. Polynomial degree bounds and
the elementary leading-term cancellation lemma give the digit-field
algebraic dichotomy. Period matrices supply algebraic existence, and
support cosets supply hidden independence.

Only the normal-form field construction, the cut property, real
closedness of $\No$, and classical real continued-fraction existence are
imported. In particular the proof of algebraic rigidity does not require
a real-closedness theorem for every Hahn field, a valuation-extension
structure theorem, a saturation theorem, or any of the repository's
advanced unrefereed results.

\subsection{The finite verification artifact}\label{ocf:sub:artifact}

The supplied program \texttt{code/verify.py} uses exact rational arithmetic,
finite generalized polynomials with rational growth exponents, and SymPy
symbolic identities. The recorded run passed \textbf{81,256 assertions}.
The main categories are as follows.

\begin{center}
\begin{tabular}{@{}lr@{}}
\toprule
Check family & Assertions\\
\midrule
Fibonacci denominator inequalities & 28,268\\
Two-step cylinder buffers & 4,294\\
Generalized-polynomial degree laws & 1,260\\
Generalized-polynomial determinant identities & 1,260\\
Rational cylinder widths & 5,694\\
Rational digit recovery & 5,694\\
\bottomrule
\end{tabular}
\end{center}

Other checks cover both endpoint-distance identities, monotonicity,
nesting, period matrices, the fractional linear difference formula, a
counterexample to multiplicativity of support projection, and the Catalan
Laurent expansion through its justified truncation order. The full
category counts, fixed seed, Python version, and SymPy version are in
\texttt{audits/verification.json}.

These are finite consistency checks of identities used by the proof. They
cannot check the proper-class existence quantifier, arbitrary infinite
supports, countable cofinality, or the absence of every polynomial
relation. Those conclusions rest on the written proofs, not on the number
of test cases.

\wtag{} In this report the program is shipped as \rpath{code/verify.py},
the recorded run as \rpath{data/verification.json} (delivered as
\texttt{audits/verification.json}) and the Makefile as
\rpath{code/Makefile}. The program writes a file only when given
\texttt{--output}; the Makefile's \texttt{check} target writes the delivery
path \texttt{audits/verification.json}. When the report was written the
program was rerun on a copy; the counts, the seed and the SymPy version
agree with the record, and only the Python version differs
(Appendix~\ref{ocf:app:suite}).

\subsection{A practical formalization route}\label{ocf:sub:formal}

A future Lean development should first isolate the ordered-field
continuant package from the surreal-specific normal-form package. The
matrix identities, cylinder inequalities, and two-step buffer need only
finite recursion, ordered-field arithmetic, and discreteness of the digit
ring. The exact-fiber proof then adds an Archimedean comparison relation
quantifying over ordinary natural numbers and the denominator-degree law.

The next genuinely surreal step is the support projection lemma, stated
as an additive subseries construction and explicitly not as a ring
homomorphism. The algebraic-error lemma should be proved separately for a
valued field, by the repeated-minimum argument in
Lemma~\ref{ocf:lem:algerror}. This keeps the main algebraic classification
independent of heavy real-closed-field automation.

No Lean file is included as though it had been checked. The repository
was not built during this task. A written proof, a finite test run, and a
kernel-checked theorem remain three different levels of evidence.

\wtag{} Of the inputs listed here, only the omnific floor of
Lemma~\ref{ocf:lem:floor} is formalized in the collection, as
\lbl{odg:thm:floor}. No continuant, cylinder, fibre or precision-ideal
statement has a Lean counterpart (Section~\ref{ocf:sec:collection}).

\section{Further research questions}
\label{ocf:sec:questions}

The following are proposed problems arising from the proved statements.
They are not all asserted to be previously published open problems. The
historical comparison is itself part of the research program.

\begin{problem}[Priority and exact equivalence of formulations]\label{ocf:q:priority}
Compare the valuation-fiber formula, full Hahn realization theorem, and
digit-field algebraic criterion against the wider literature following
\cite{Roggeman}. Which statements are genuinely new, and which are hidden
consequences of earlier associated-class or generalized-fraction
formulations? A useful outcome would be a theorem-by-theorem concordance,
not merely a list of citations.
\end{problem}

\begin{problem}[Proper subfields of a full Hahn field]\label{ocf:q:subfields}
Characterize subfields $L\subsetneq K_G$ containing $\R$ and closed under
omnific floor for which every admissible code in $L\cap\Oz$ has a
realization in $L$. Full Hahn fields have the property by
Theorem~\ref{ocf:thm:hahn}; arbitrary subfields do not follow from its proof.
Determine which support-completeness or truncation assumptions suffice.
\end{problem}

\begin{problem}[Algebraic existence for a nonperiodic code]\label{ocf:q:nonperiodic}
Theorem~\ref{ocf:thm:algebraic} determines the size of the algebraic fiber
once its nonemptiness is known. Find checkable conditions on nonperiodic
omnific digits that imply, or rule out, algebraic existence over the
countable digit field. In particular, investigate codes with cofinal
degrees and a finitely described recurrence among their digits.
\end{problem}

\begin{problem}[A correctly scoped Lagrange converse]\label{ocf:q:lagrange}
For a specified full Hahn field and a specified coefficient subfield,
which algebraic points of degree two have eventually periodic omnific
floor itineraries? The coefficient field and the allowed digits must be
fixed in advance. Theorem~\ref{ocf:thm:periodroot} proves an existence direction
for a given periodic code, not this converse.
\end{problem}

\begin{problem}[Simplicity versus dynamical selection]\label{ocf:q:simplicity}
When does the simplest cylinder separator $\{L\mid R\}$ coincide with the
period-fixed algebraic representative of an eventually periodic code?
How does this comparison behave under finite prefix maps? Establishing
precise hypotheses would connect the present valuation classification
with the simplicity-based selections in earlier work. Roggeman's residue
and periodicity theorems must be compared before any instance is
classified as a previously unresolved problem.
\end{problem}

\begin{problem}[Iteration beyond the first block of length $\omega$]\label{ocf:q:transfinite}
Design and compare limit-stage remainder rules that refine a full fiber
without discarding the original point. The transfinite problem is not to
redeclare the first-block cut to be its unique value. It is to retain and
encode an element of the precision ideal, compatibly with earlier
generalized continued-fraction constructions.
\end{problem}

\begin{problem}[Other integer parts]\label{ocf:q:integerparts}
Which parts of the classification remain true for a different discretely
ordered integer part of an ordered field? The finite cylinder and buffer
lemmas are robust, but the degree identity, the ordinary-tail argument,
and the support projection use special features of $\Oz$. An axiomatic
characterization of exactly those features would be valuable.
\end{problem}

\begin{problem}[Primary precision ideals]\label{ocf:q:primary}
Theorem~\ref{ocf:thm:prime} and Corollary~\ref{ocf:cor:radical}\footnote{\wtag{}
The source wrote ``Theorems'' for this pair; the second is a corollary.}
classify primality and the
radical directly from the degree cut. Characterize when a precision ideal
is primary in its valuation ring, and describe the torsion and
annihilators in the associated error quotient at higher value-group rank.
The rank-one examples suggest useful test cases but do not decide the
higher-rank problem.
\end{problem}

\begin{problem}[A genuinely surcomplex continued-fraction algorithm]\label{ocf:q:surcomplex}
Choose an explicit nearest-Gaussian-omnific-integer rule and determine
whether any analogue of the exact fiber theorem survives. A central
obstacle is that complex denominators can have leading cancellation, so
positivity cannot supply either monotonicity or the degree recurrence.
The coordinatewise result of Section~\ref{ocf:sec:complex} does not solve this.
\end{problem}

\begin{problem}[Effective precision and representations]\label{ocf:q:effective}
For computably presented exponent groups and exact digit notations,
determine the complexity of finding an index at which a specified error
becomes visible. Separate the problem of deciding equality of digits from
the cofinality problem for their cumulative degrees. A finite program for
rational exponents is not an oracle for arbitrary surreal descriptions.
\end{problem}

\begin{problem}[Birthday bounds]\label{ocf:q:birthday}
Bound the birthday of a simplest realization from the birthdays of its
digits and their support cuts. Compare those bounds with the birthdays
of the explicit quadratic representative in the periodic case. Ordinary
field degree and surreal birthday should not be identified, and no
general product-birthday conjecture is resolved by the present work.
\end{problem}

\begin{problem}[Formal proof and reusable interfaces]\label{ocf:q:formal}
Formalize the buffer lemma, exact fiber theorem, support-projection
existence argument, and algebraic dichotomy in the repository's actual
surreal field. The decisive interface question is how to expose
normal-form subseries and their leading valuations without introducing
unproved summability or class-size assumptions.
\end{problem}

\wtag{} All twelve questions are open in the collection
(Section~\ref{ocf:sec:collection}). For Question~\ref{ocf:q:formal}, the
omnific floor is already formalized (\lbl{odg:thm:floor}); none of the
statements the question names is.

\section{Conclusion}\label{ocf:sec:conclusion}

An omnific continued fraction is most naturally treated as precision data.
The exact data are the coset $x+I_{\codea}$, with an explicit lower cut on
allowed error valuations. From that viewpoint three seemingly conflicting
phenomena become compatible: every full surreal infinite fiber is large;
a full Hahn workspace can make a code unique; and a countable digit field
can single out one algebraic member of a much larger fiber.

The two most useful thresholds are different. General uniqueness in a
full Hahn field detects countable cofinality of its value group, whereas
periodic uniqueness detects an order unit. Algebraic uniqueness over the
digit field is controlled instead by self-doubling of the cumulative
degrees and visibility of the initial digit. These precise distinctions,
rather than the already known failure of ordinary uniqueness, are the
substantive outcome of this investigation.

\section{Consolidated ledger of non-claims}\label{ocf:sec:nonclaims}

\wtag{} The source states its limitations where they arise, and each is
kept in place. This ledger collects them with their locations.
``README'' is the delivered README, which is not shipped (its content is in
the collection README); \rpath{proof\_and\_scope.md} and
\rpath{source\_audit.md} are the delivered audits, shipped unchanged.
Items S1--S26 are the source's; items W1--W8 were added when the report
joined the collection.
\begin{enumerate}[label=S\arabic*.,leftmargin=2.6em,itemsep=2pt]
\item Basic nonuniqueness of surreal continued fractions and proper-class
fibres are \emph{not} new. Roggeman's paper \cite{Roggeman} already treats
the omnific floor algorithm, associated classes, their size and structure,
periodic fractions and a generalized continuation (abstract; title page;
Section~\ref{ocf:sub:known}; after Corollary~\ref{ocf:cor:neverunique};
Section~\ref{ocf:sec:examples}; README; \rpath{source\_audit.md}).
\item The novelty and priority of the proposed contributions (the precision
ideal, the full-Hahn realization theorem, the prime criterion and the
algebraic dichotomy) are not independently certified. No literature search
is claimed to establish them; a wider concordance could show some of them
implicit in earlier work; no label such as ``breakthrough'' is claimed
(title page; Section~\ref{ocf:sec:intro}; Question~\ref{ocf:q:priority};
README; \rpath{source\_audit.md}; \rpath{proof\_and\_scope.md}).
\item No certified resolution of a named longstanding open problem is
claimed. The research questions are proposed, not all previously
published open problems (title page; Section~\ref{ocf:sec:questions};
Appendix~\ref{ocf:app:notation}; README; \rpath{source\_audit.md}).
\item The finite tests are not a Lean verification and not a proof of the
infinite-support or class-sized assertions. They cannot check the
proper-class existence quantifier, arbitrary infinite supports, countable
cofinality, the absence of every polynomial relation, the complete proof or
novelty (title page; Section~\ref{ocf:sub:artifact};
Appendix~\ref{ocf:app:notation}; README; \rpath{proof\_and\_scope.md};
the \texttt{not\_certified} list of \rpath{data/verification.json}).
\item No Lean file is included as checked, the repository was not built,
and no proof-assistant certificate exists. A written proof, a finite test
run and a kernel-checked theorem are three different levels of evidence
(Section~\ref{ocf:sub:formal}; README; \rpath{proof\_and\_scope.md}).
\item There has been no independent referee review; the delivered proof audit
is a self-audit (README; \rpath{proof\_and\_scope.md};
\rpath{source\_audit.md}).
\item The sequence length is exactly $\omega$ and every index is an ordinary
natural number. No limit-stage continuation of the algorithm is included or
claimed (Section~\ref{ocf:sec:intro}; Question~\ref{ocf:q:transfinite};
README).
\item No nearest-Gaussian-integer algorithm is claimed. Section~\ref{ocf:sec:complex}
is a coordinatewise theory, not a rotation-invariant coding of all
surcomplex points (Section~\ref{ocf:sec:intro}; after
Theorem~\ref{ocf:thm:complex}; Question~\ref{ocf:q:surcomplex}; README).
\item Realization is proved in \emph{full} Hahn fields only. It is not a
completeness assertion about arbitrary non-Archimedean ordered fields, and
arbitrary subfields of $K_G$ do not follow from the proof
(Section~\ref{ocf:sec:intro}; Section~\ref{ocf:sec:hahn};
Question~\ref{ocf:q:subfields}; README; \rpath{proof\_and\_scope.md}).
\item The support projection $\pi_G$ is additive, not multiplicative and not
a field map; the proof never projects a product or a reciprocal
(Remark~\ref{ocf:rem:projection}; \rpath{proof\_and\_scope.md}).
\item Algebraic existence is not asserted for nonperiodic codes; for an
arbitrary code $A_{\codea}$ is not known to be nonempty
(Section~\ref{ocf:sec:algebraic}; Theorem~\ref{ocf:thm:algebraic};
Question~\ref{ocf:q:nonperiodic}; README; \rpath{proof\_and\_scope.md}).
\item No Lagrange converse is claimed: it is not claimed that every quadratic
surreal has an eventually periodic omnific itinerary
(Question~\ref{ocf:q:lagrange}; README; \rpath{proof\_and\_scope.md}).
\item It is not claimed that every point with an eventually periodic
itinerary is quadratic: periodic digits do not force complete quotients to
recur (after Theorem~\ref{ocf:thm:periodroot}; \rpath{proof\_and\_scope.md}).
\item The simplest separator $\{L\mid R\}$ of Corollary~\ref{ocf:cor:realization}
is not claimed to be the period-fixed root, nor an order-topological limit
(after Corollary~\ref{ocf:cor:realization}; Question~\ref{ocf:q:simplicity};
\rpath{proof\_and\_scope.md}).
\item Unrestricted surreal fibres are never unique
(Corollary~\ref{ocf:cor:neverunique}); no unique unrestricted realization
is claimed (\rpath{proof\_and\_scope.md}).
\item Countable cofinality is not claimed to imply an order unit; the example of
Section~\ref{ocf:sec:examples} separates them (after
Theorem~\ref{ocf:thm:orderunit}; \rpath{proof\_and\_scope.md}).
\item Theorem~\ref{ocf:thm:independence} gives, for each prescribed set
field and set cardinal separately, one independent family. It does not give
a proper-class transcendence basis, a family independent over all set
subfields at once, or a first-order definable or computable family
(Section~\ref{ocf:sec:prelim}; Remark~\ref{ocf:rem:independence};
\rpath{proof\_and\_scope.md}).
\item Every support, sequence, generated field and cut used in a
construction is a set. No large-cardinal assumption, class-length
recursion for the algorithm, or proper-class choice is used; quotient rings
are taken in set-sized Hahn workspaces, and no object of proper-class
cosets is assumed (Section~\ref{ocf:sec:prelim}; after
Theorem~\ref{ocf:thm:prime}; README; \rpath{proof\_and\_scope.md}).
\item Proposition~\ref{ocf:prop:topology} is an elementary class-size
obstruction, not a new topology theorem (after
Proposition~\ref{ocf:prop:topology}).
\item The continuant identities, the real periodicity theorem, the periodic
fixed-point mechanism and the radical argument are classical; they are
rederived because the digits are infinite, not claimed as new
(Section~\ref{ocf:sub:known}; after Theorem~\ref{ocf:thm:periodroot};
proof of Corollary~\ref{ocf:cor:radical}).
\item Uniqueness always names its ambient field and algebraicity its
coefficient field; the algebraic fibre is relative to the digit field, not
to the whole surreal field; prime-ideal claims refer to the finite
valuation ring, not to the field (Section~\ref{ocf:sec:algebraic};
Appendix~\ref{ocf:app:notation}; README; \rpath{proof\_and\_scope.md}).
\item Ordinary real coefficient convergence is not claimed to imply order
convergence in $K_G$ when $G\ne\{0\}$ (after
Corollary~\ref{ocf:cor:convergence}).
\item Proposition~\ref{ocf:prop:mobius} does not claim that a fractional linear
map preserves digit codes, and its linear threshold law needs
\eqref{ocf:eq:smalltransport}. Equal coordinate ideals give a possibly
nonprincipal intersection of valuation balls, not one ball of realized
radius (after Theorem~\ref{ocf:thm:complex}; after
Proposition~\ref{ocf:prop:mobius}; \rpath{proof\_and\_scope.md}).
\item Field degree and birthday are not identified, and no general
product-birthday conjecture is resolved (Question~\ref{ocf:q:birthday}).
\item Only normal forms, the cut property, real closedness of $\No$ and real
continued-fraction existence are imported; no unreviewed advanced
repository result is used, and the repository comparison was a focused
catalogue comparison, not a declaration-by-declaration or whole-history
audit (Section~\ref{ocf:sub:known}; Section~\ref{ocf:sec:verification};
README; \rpath{source\_audit.md}).
\item Bibliographic caveats: the journal metadata of \cite{Roggeman} were
not established and are not invented; Eberl's formalization \cite{Eberl}
concerns real continued fractions and certifies nothing here; the
foundations \cite{Conway,Gonshor} are imported, not reproved. The delivered
build and layout record (\rpath{data/build\_and\_layout.json}) is an
artifact check, not a proof or priority certification (bibliography;
\rpath{source\_audit.md}; \rpath{data/build\_and\_layout.json}).
\end{enumerate}
\begin{enumerate}[label=W\arabic*.,leftmargin=2.6em,itemsep=2pt]
\item No \texttt{ocf:} label has a Lean mapping in \rpath{FORMALIZATION.md}.
Independent proof review and formalization are pending. The only
formalized ingredient is the omnific floor, through \lbl{odg:thm:floor}
(Section~\ref{ocf:sec:collection}).
\item Lemma~\ref{ocf:lem:floor} is not new to the collection: it is
\lbl{odg:thm:floor}, \lbl{dsn:prop:floor}, \lbl{cas:eq-floor} and
\lbl{onot:eq:floor}. It is printed once (Section~\ref{ocf:sec:collection}).
\item The support-coset argument of Theorem~\ref{ocf:thm:independence} is
credited to the independent-copies and bounded-support reports; no
priority is claimed for it relative to them. The comparison with the
expanding-polynomial report's itinerary fibres and order-unit threshold is
a parallel of shape: neither result is derived from the other, and no
priority is claimed either way (Section~\ref{ocf:sec:collection}).
\item No named question of any report is answered; the twelve questions of
Section~\ref{ocf:sec:questions} stay open (Section~\ref{ocf:sec:collection}).
\item The source's statements about the collection are those of its pin
\texttt{958b5c4}; they were rechecked and remain true at the placement
commit, and are kept as provenance (Appendix~\ref{ocf:app:pinned}).
\item The proofs were read when the report was written
(Appendix~\ref{ocf:app:checked}). That reading is not an independent proof
review.
\item The rerun of the finite suite (Appendix~\ref{ocf:app:suite}) agrees
with the record except for the Python version. A passing rerun certifies
the finite identities only, not any theorem of the report.
\item Roggeman's paper was not read for this report. The attributions to
its Sections 3A, 3H and 3I and the statement that its scan was inspected
are the source's (Appendix~\ref{ocf:app:literature}).
\end{enumerate}

\appendix
\section{Notation and scope checklist}\label{ocf:app:notation}

\begin{longtable}{@{}p{0.22\textwidth}p{0.73\textwidth}@{}}
\toprule
Notation & Meaning\\
\midrule
\endhead
$\No$, $\Oz$, $\SC$ & Surreal field, omnific integer ring, and $\No[i]$.\\
$\dgr x$, $v(x)$ & Greatest growth exponent and its negative; $v(0)=+\infty$.\\
$\Ofin$, $\mm$ & Finite surreals and infinitesimals (\wtag{} $\Ofin$ is the source's $\mathcal O$).\\
$K_G$ & Full real-coefficient Hahn field with growth exponents in the set group $G$.\\
$\codea$, $\F(\codea)$ & An ordinary $\omega$-indexed admissible code and all its nonterminating realizations.\\
$p_n,q_n$ & Finite continuants; $q_0=1$.\\
$r_n,s_n,C_n$ & Convergent, other cylinder endpoint, and open prefix cylinder.\\
$d_n$ & $\dgr q_n=\sum_{j=1}^n\dgr a_j$; the first digit does not contribute.\\
$I_{\codea}$ & Errors whose valuations exceed $2d_n$ for every ordinary $n$.\\
$K_{\codea}$ & Countable field generated over $\Q$ by the digits themselves.\\
$D_0$ & Degree of the first digit, with value zero assigned if that digit is zero.\\
$E,D_*$ & Total degree in a period and largest degree among all displayed digits.\\
\bottomrule
\end{longtable}

All assertions about uniqueness name their ambient field. All assertions
about algebraicity name their coefficient field. All arbitrary-support
series are Hahn normal forms, not analytic limits. Prime-ideal claims
refer to the finite valuation ring, not to the whole field. Finite tests
are not formal verification. The proposed research questions are not
advertised as newly solved published conjectures.

\section{Provenance of this report}\label{ocf:app:report}

\subsection{One manuscript}\label{ocf:app:onesource}

\wtag{} This report is built from \emph{one} manuscript: manuscript~09 of
batch~33 of the collection, the archive
\texttt{surreal\_continued\_fractions\_research} (inner directory
\texttt{surreal\_cf\_article}), titled ``Exact Precision and Algebraic
Rigidity of Omnific Continued Fractions'', a 22-page PDF dated
23~September 2026 and pinned to repository commit
\nolinkurl{958b5c4865819bd55ea1f5ffc050282aea7ef570}. The archive was
delivered in commit \texttt{73043eb} and placed in \texttt{aa9c891} as a
new report. There was no second source: nothing was merged and nothing was
selected away, and every result, proof, example, question and limitation of
the manuscript is printed here.

The manuscript contributed everything in
Sections~\ref{ocf:sec:intro}--\ref{ocf:sec:conclusion} and
Appendix~\ref{ocf:app:notation} except the text marked \wtag: the exact
floor, continuant cylinders and their widths
(Lemma~\ref{ocf:lem:floor} to Corollary~\ref{ocf:cor:realization}); the
two-step buffer, the exact convergent error and the exact fibre
$\F(\codea)=x+I_{\codea}$ (Lemmas~\ref{ocf:lem:buffer},
\ref{ocf:lem:degrees}, Theorem~\ref{ocf:thm:fiber},
Corollary~\ref{ocf:cor:neverunique}, Proposition~\ref{ocf:prop:topology});
realization in every full Hahn field with its uniqueness criterion and the
countable-cofinality threshold (Theorem~\ref{ocf:thm:hahn},
Corollaries~\ref{ocf:cor:cofinality} and~\ref{ocf:cor:convergence}); the
prime criterion (Proposition~\ref{ocf:prop:prescribed},
Theorem~\ref{ocf:thm:prime}, Corollary~\ref{ocf:cor:radical}); algebraic
rigidity over the digit field (Lemmas~\ref{ocf:lem:algerror},
\ref{ocf:lem:digitfield}, Theorem~\ref{ocf:thm:algebraic}); periodic codes
(Theorems~\ref{ocf:thm:periodroot}, \ref{ocf:thm:periodiccriterion},
\ref{ocf:thm:orderunit}, Corollary~\ref{ocf:cor:pureperiod}); six worked
examples; hidden independence (Theorem~\ref{ocf:thm:independence}); the
coordinatewise surcomplex fibre and precision transport
(Theorem~\ref{ocf:thm:complex}, Corollary~\ref{ocf:cor:complexindependence},
Proposition~\ref{ocf:prop:mobius}); and twelve research questions.

The package held the article, its PDF, a README, a Makefile, a
requirements file, \texttt{code/verify.py} and an \texttt{audits/}
directory with two JSON records and two markdown audits. It had no checksum
manifest. The article was renamed \rpath{article.tex} and rewritten as
described below, the README was replaced by the collection README, and the
delivered PDF is not shipped: \rpath{article.pdf} is a build of this text.
The other files are shipped byte-identical, under the paths of
Appendix~\ref{ocf:app:files}.

The editorial changes are these.
\begin{enumerate}[label=(\arabic*)]
\item Every label carries the prefix \texttt{ocf:}. The source's 77 labels
are kept, unchanged after the prefix. Thirty-five were added: the new
sections and subsections \lbl{ocf:sub:known}, \lbl{ocf:sec:conventions},
\lbl{ocf:sec:collection}, \lbl{ocf:sec:nonclaims}, \lbl{ocf:sec:conclusion},
\lbl{ocf:sub:artifact}, \lbl{ocf:sub:formal}, \lbl{ocf:app:notation},
\lbl{ocf:app:report}, \lbl{ocf:app:onesource}, \lbl{ocf:app:pinned},
\lbl{ocf:app:literature}, \lbl{ocf:app:checked}, \lbl{ocf:app:suite} and
\lbl{ocf:app:files}; the unlabelled environments \lbl{ocf:def:Oz},
\lbl{ocf:def:digits}, \lbl{ocf:def:ideal}, \lbl{ocf:rem:projection},
\lbl{ocf:rem:twocases}, \lbl{ocf:cor:convergence},
\lbl{ocf:rem:independence} and \lbl{ocf:cor:complexindependence}; and the
twelve questions \lbl{ocf:q:priority}, \lbl{ocf:q:subfields},
\lbl{ocf:q:nonperiodic}, \lbl{ocf:q:lagrange}, \lbl{ocf:q:simplicity},
\lbl{ocf:q:transfinite}, \lbl{ocf:q:integerparts}, \lbl{ocf:q:primary},
\lbl{ocf:q:surcomplex}, \lbl{ocf:q:effective}, \lbl{ocf:q:birthday} and
\lbl{ocf:q:formal}. Labels on existing environments change no number.
\item One symbol was renamed: the source's $\mathcal O$ (finite surreals) is
$\Ofin$, and its $\mathcal O_\C$ is $\Ofin[i]$
(Section~\ref{ocf:sec:conventions}).
\item One mistyped cross-reference was corrected: the source's
``Theorems 6.2 and 6.3'' in Question~\ref{ocf:q:primary} names a theorem
and a corollary (footnote there).
\item The title page names the report as a single-source report of the
collection. Sections~\ref{ocf:sec:conventions}, \ref{ocf:sec:collection}
and~\ref{ocf:sec:nonclaims} and this appendix are new. Sentences marked
\wtag{} were added at the end of Section~\ref{ocf:sub:known}; after
Lemma~\ref{ocf:lem:floor}; after Theorem~\ref{ocf:thm:orderunit}; after
Remark~\ref{ocf:rem:independence}; in
Sections~\ref{ocf:sub:artifact} and~\ref{ocf:sub:formal}; after
Question~\ref{ocf:q:formal}; and in the notation table of
Appendix~\ref{ocf:app:notation}.
\end{enumerate}
No mathematical statement was changed. Where the edit had to choose: the
manuscript was placed as a new report, not as an addition. Its floor
lemma is the Diophantine report's, but that is one imported ingredient,
and the expanding-polynomial report shares the fibre-plus-ideal shape and
the order-unit threshold for a different coding; neither report names a
question that this manuscript answers. It was placed among the surreal
reports because its fields are subfields of $\No$; the coordinatewise
surcomplex section is a supplement. The title is kept.

\subsection{Repository statements at the pin}\label{ocf:app:pinned}

The source compared itself with the root README, the catalogue and two
surcomplex report READMEs at its pin (\rpath{source\_audit.md},
\cite{SurrealRepo}). The pin is an ancestor of the placement commit
\texttt{aa9c891}, which is 71 commits later. Checked there:
\begin{itemize}[leftmargin=*,itemsep=3pt]
\item No report treats surreal or omnific continued fractions, and no report
cites Roggeman; the only other continued fractions and continuants are
those listed in Section~\ref{ocf:sec:collection}. The source's statement
that no continued-fraction report exists stands.
\item The omnific floor, degree and normal-form infrastructure it mentions
exists; the floor is formalized (\lbl{odg:thm:floor}).
\item No report supplies Theorems~\ref{ocf:thm:fiber}, \ref{ocf:thm:hahn},
\ref{ocf:thm:prime}, \ref{ocf:thm:algebraic}, \ref{ocf:thm:periodiccriterion}
or~\ref{ocf:thm:orderunit}. No stale statement of the source needed
correction.
\end{itemize}

\subsection{References checked for this report}\label{ocf:app:literature}

\begin{itemize}[leftmargin=*,itemsep=3pt]
\item \emph{Roggeman} \cite{Roggeman}. A web search finds the
author-hosted record and PDF at the address given, under the stated title
and author. The paper was not read for this report, and its journal was
not established; the attributions to its sections are the source's.
\item Conway \cite{Conway}, Gonshor \cite{Gonshor}, Krishnan
\cite{Krishnan} and Eberl \cite{Eberl} were not checked.
\end{itemize}

\subsection{Proofs read when the report was written}\label{ocf:app:checked}

Every proof of Sections~\ref{ocf:sec:prelim}--\ref{ocf:sec:complex} was read
line by line when the report was written. No error was found. Rechecked in
detail: the endpoint distances \eqref{ocf:eq:distR}--\eqref{ocf:eq:distS}
and the identity $\lvert x-s_n\rvert=1/(q_{n+1}(q_{n+1}u+q_n))$ in the case
$a_{n+1}=1$ of Lemma~\ref{ocf:lem:buffer}; the degree induction of
Lemma~\ref{ocf:lem:degrees}; the bound $2d_n\le2(r-1)D_*+2(k+1)E<(2r+2)D_*$
behind $\delta=b^{-(2r+2)}\in I_{\codea}$ in
Theorem~\ref{ocf:thm:periodiccriterion}; the Catalan coefficients of
\eqref{ocf:eq:rho} against the binomial series of
$\omega\sqrt{1+4\omega^{-2}}$; the rank-one boundary
\eqref{ocf:eq:rankoneboundary}; and the degrees
$d_n=\sum_{j\le n}\omega^j$ of the example without an order unit. The only
defect found was the mistyped reference corrected in
Question~\ref{ocf:q:primary}. This reading is not an independent proof
review.

\subsection{The verification suite}\label{ocf:app:suite}

\rpath{code/verify.py} (Python~3.10 or later, SymPy; the record used
Python~3.13.5 and SymPy~1.14.0, pinned in \rpath{data/requirements.txt})
checks finite rational continued fractions, finite generalized polynomials
with rational exponents and SymPy identities, with the fixed seed
20260923. It raises on the first failed check, prints its JSON report, and
writes a file only with \texttt{--output}. It was rerun on a copy for this
report under Python~3.14.4 and SymPy~1.14.0: exit code~0,
\texttt{"status": "passed"}, 81,256 assertions in the same 20 categories
with the same counts as \rpath{data/verification.json}; the only difference
is the recorded Python version.

\subsection{Files}\label{ocf:app:files}

\begin{center}\small
\begin{tabular}{@{}ll@{}}
\toprule
File & Delivered as\\
\midrule
\rpath{article.tex} & \texttt{article.tex}, rewritten as this report\\
\rpath{article.pdf} & a build of this text, not the delivered PDF\\
\rpath{README.md} & the collection README, replacing the delivered one\\
\rpath{proof\_and\_scope.md} & \texttt{audits/proof\_and\_scope.md}\\
\rpath{source\_audit.md} & \texttt{audits/source\_audit.md}\\
\rpath{code/verify.py} & \texttt{code/verify.py}\\
\rpath{code/Makefile} & \texttt{Makefile}\\
\rpath{data/requirements.txt} & \texttt{requirements.txt}\\
\rpath{data/verification.json} & \texttt{audits/verification.json}\\
\rpath{data/build\_and\_layout.json} & \texttt{audits/build\_and\_layout.json}\\
\bottomrule
\end{tabular}
\end{center}

The shipped delivered files still use the delivery paths: the docstring
of \rpath{code/verify.py} and the Makefile's \texttt{check} target write
\texttt{audits/verification.json}, the Makefile expects to run from the
report root, and \rpath{proof\_and\_scope.md} names
\texttt{verification.json} without a directory.
\rpath{data/build\_and\_layout.json} records the delivered 22-page build,
not this one.

\clearpage
\phantomsection
\addcontentsline{toc}{section}{References}
\begin{thebibliography}{9}

\bibitem{Conway}
J.~H. Conway, \emph{On Numbers and Games}, Academic Press, 1976;
second edition, A K Peters, 2001.
Standard source for the surreal construction, normal forms, and field structure.

\bibitem{Gonshor}
H. Gonshor, \emph{An Introduction to the Theory of Surreal Numbers},
London Mathematical Society Lecture Note Series 110, Cambridge University
Press, 1986.

\bibitem{Roggeman}
Y. Roggeman, \emph{Continued fractions and surreal numbers}, 1985,
scanned article, pp.~31--70; received 15 June 1985.
Author-supplied full text available from
\href{https://www.researchgate.net/publication/268854423_Continued_fractions_and_surreal_numbers}{the author's ResearchGate record}.
The complete 21-page scan (two printed pages per scan after the first)
was inspected; journal metadata was not needed for the comparison.

\bibitem{Krishnan}
Gautam Gopal Krishnan, \emph{Continued Fractions}, Cornell University notes,
22 August 2016.
\url{https://pi.math.cornell.edu/~gautam/ContinuedFractions.pdf}.

\bibitem{Eberl}
M. Eberl, \emph{Continued Fractions}, Archive of Formal Proofs,
20 March 2024.
\url{https://isa-afp.org/entries/Continued_Fractions.html}.
The cited formalization concerns real continued fractions, not the
surreal theorems in this article.

\bibitem{SurrealRepo}
V. Reshetnikov, \emph{Surreal}, research repository,
\url{https://github.com/VladimirReshetnikov/Surreal},
snapshot \texttt{958b5c4865819bd55ea1f5ffc050282aea7ef570}, inspected
23 September 2026. Focused sources: root \texttt{README.md},
\texttt{docs/README.md}, and selected report READMEs.

\end{thebibliography}
\end{document}
